\chapter[Operatory nieograniczone i teoria spektralna]{Operatory nieograniczone, samosprzężoność, spektrum i półgrupy}
\label{chap:pde12}

\providecommand{\atlasref}[2]{#2}

W rozdziale~\ref{chap:pde-variational} zagadnienie Dirichleta rozwiązaliśmy
w przestrzeni energetycznej $H_0^1(\Omega)$. Teraz zapytamy, kiedy
formalny operator $-\Delta$ jest operatorem samosprzężonym w
$L^2(\Omega)$, jak wyglądają jego wartości własne i dlaczego
jego widmo steruje ewolucją cieplną i falową. Dalsze rachunki
zostaną zapisane najpierw dla rzeczywistych przestrzeni Hilberta.
W argumentach z operatorami unitarnymi i widmem zespolonym
korzystamy z ich zespolenia. W zespolonej przestrzeni Hilberta
iloczyn skalarny przyjmujemy liniowy w pierwszym argumencie.

\section{Dziedzina jest częścią operatora}
\label{sec:pde12-domain}

\begin{definicja}[Operator liniowy, domknięty i domykalny]
\label{def:pde12-operator}
Niech $H$ będzie przestrzenią Hilberta. Operatorem liniowym
$T$ w $H$ nazywamy parę składającą się z podprzestrzeni
$D(T)\subset H$, zwanej \emph{dziedziną}, oraz odwzorowania
liniowego $T:D(T)\to H$. Wykresem $T$ nazywamy
$\mathcal G(T)=\{(u,Tu):u\in D(T)\}\subset H\times H$.
Wykres jest podprzestrzenią liniową $H\times H$:
dla $u,v\in D(T)$ i skalarów $a,b$ mamy
$a(u,Tu)+b(v,Tv)=(au+bv,T(au+bv))$.
Operator $T$ jest \emph{domknięty}, jeżeli $\mathcal G(T)$ jest
domknięty w $H\times H$. Jest \emph{domykalny}, jeżeli
domknięcie $\mathcal G(T)$ jest wykresem pewnego operatora;
operator ten oznaczamy przez $\overline T$.
Równoważnie $T$ jest domykalny wtedy i tylko wtedy, gdy
dla każdego ciągu $(u_n)$ z $D(T)$ warunki
$u_n\to0$ i $Tu_n\to v$ implikują $v=0$.
Istotnie, domknięcie wykresu jest podprzestrzenią liniową
i jest wykresem dokładnie wtedy, gdy nad pierwszą współrzędną
$0$ zawiera jedynie parę $(0,0)$: dwie pary $(u,v_1)$
i $(u,v_2)$ dałyby przez odjęcie $(0,v_1-v_2)$.
Odwrotnie, jeśli warunek ten zachodzi, na rzucie domknięcia
wykresu na pierwszą współrzędną (podprzestrzeni liniowej)
możemy jednoznacznie przypisać każdemu $u$ drugą współrzędną.
Liniowość domknięcia wykresu daje liniowość tego przypisania.
Ponieważ $H\times H$ jest przestrzenią metryczną,
$(0,v)$ należy do domknięcia wykresu dokładnie wtedy,
gdy jest granicą pewnego ciągu $(u_n,Tu_n)$.
W przypadku domykalności operator $\overline T$ jest
liniowy i domknięty, rozszerza $T$ i zawiera się w każdym
domkniętym rozszerzeniu $T$, ponieważ wykres takiego
rozszerzenia zawiera domknięcie wykresu $T$.
\end{definicja}

\begin{lemat}[Kryterium domkniętości i norma grafowa]
\label{lem:pde12-graph}
Dla operatora liniowego $T:D(T)\subset H\to H$
następujące warunki są równoważne:
(i) $T$ jest domknięty; (ii) ilekroć $u_n\in D(T)$,
$u_n\to u$ i $Tu_n\to v$ w $H$, to $u\in D(T)$ oraz $Tu=v$;
(iii) przestrzeń $D(T)$ z normą
$\|u\|_{D(T)}=(\|u\|_H^2+\|Tu\|_H^2)^{1/2}$ jest zupełna.
\end{lemat}
\begin{proof}
Zbieżność $(u_n,Tu_n)\to(u,v)$ w iloczynie $H\times H$
jest równoważna dwóm zbieżnościom z (ii), więc (i) i (ii)
są równoważne. Odwzorowanie $u\mapsto(u,Tu)$ jest izometrią
$D(T)$ z normą grafową na $\mathcal G(T)$ z normą iloczynową:
jest liniowe, różnowartościowe dzięki pierwszej współrzędnej
i na mocy definicji wykresu surjektywne. Norma grafowa
pochodzi więc od iloczynu skalarnego
$(u,v)_{D(T)}=(u,v)_H+(Tu,Tv)_H$.
Domknięta podprzestrzeń przestrzeni zupełnej jest zupełna.
Odwrotnie, jeśli $(u_n,Tu_n)$ zbiega w $H\times H$,
to $(u_n)$ jest ciągiem Cauchy'ego w normie grafowej.
Z zupełności $D(T)$ ma granicę $w\in D(T)$ w tej normie,
więc $(u_n,Tu_n)\to(w,Tw)$ w $H\times H$.
Jednoznaczność granicy daje przynależność pierwotnej
granicy do wykresu, który jest zatem domknięty. To daje (i)$\Leftrightarrow$(iii).
\end{proof}

\begin{przyklad}[Pochodna i dwie różne dziedziny]
\label{ex:pde12-derivative}
W $H=L^2(0,1;\C)$ rozważmy operator $D_0u=-iu'$ z dziedziną
$D(D_0)=H_0^1(0,1;\C)$. Dziedzina jest gęsta, ponieważ
zawiera $C_c^\infty(0,1;\C)$. Norma grafowa jest dokładnie
normą $H^1$, gdyż $\|-iu'\|_2=\|u'\|_2$.
Przestrzeń $H_0^1$ jest domknięta w zupełnej przestrzeni
$H^1$ (twierdzenie~\ref{thm:pde9-banach}), więc
lemat~\ref{lem:pde12-graph} daje domkniętość $D_0$.
Operator jest też symetryczny:
dla testów $u,v$ całkowanie przez części daje
\[
 (D_0u,v)=-i\int_0^1u'\overline v\,\mathrm d\lambda_1
 =i\int_0^1u\overline{v'}\,\mathrm d\lambda_1
 =(u,D_0v).
\]
Obie skrajne formy są ciągłe na $H^1\times H^1$
z nierówności Cauchy'ego--Schwarza, więc przez gęstość
testów w $H_0^1$ równość zachodzi dla całej dziedziny.
Nie jest jednak samosprzężony. W istocie
$D(D_0^*)=H^1(0,1;\C)$ i $D_0^*v=-iv'$.
Aby wykazać to ostatnie, dla $v\in D(D_0^*)$ podstawiamy
najpierw $u\in C_c^\infty(0,1)$: istnienie $g\in L^2$
takiego, że $\int_0^1(-iu')\overline v\,\mathrm d\lambda_1
=\int_0^1u\overline g\,\mathrm d\lambda_1$ oznacza
$i(\overline v)'=\overline g$ dystrybucyjnie: z definicji
pochodnej dystrybucji lewa strona działa na test $u$ przez
$-i\int_0^1u'\overline v\,\mathrm d\lambda_1$.
Po sprzężeniu mamy $-iv'=g\in L^2$, a więc $v\in H^1$.
Odwrotnie, dla $v\in H^1$ definicja słabej pochodnej daje
$(D_0u,v)=(u,-iv')$ najpierw dla testów $u$.
Dla $u\in H_0^1$ wybieramy testy $u_n\to u$ w $H^1$.
Różnice obu stron ograniczają odpowiednio
$\|u_n'-u'\|_2\|v\|_2$ i $\|u_n-u\|_2\|v'\|_2$,
więc tożsamość przechodzi na całą dziedzinę $D_0$.
To dowodzi obu inkluzji dziedzin i wzoru na sprzężenie.
W szczególności stała funkcja $1$ należy do $D(D_0^*)$,
ale nie do $D(D_0)$, gdyż jej ślad na końcach odcinka
nie jest zerowy (twierdzenie~\ref{thm:pde10-trace-zero}).
\end{przyklad}

Ogólne kryterium samosprzężoności gęsto określonych operatorów symetrycznych przez surjektywność operatorów przesuniętych podano w \atlasref{FA-33::thm:selfadjoint-range33}{FA-33: kryterium samosprzężoności}.

\begin{definicja}[Operator sprzężony i samosprzężoność]
\label{def:pde12-adjoint}
Niech $T:D(T)\subset H\to H$ ma gęstą dziedzinę.
Element $v\in H$ należy do $D(T^*)$ wtedy i tylko wtedy,
gdy istnieje $g\in H$ takie, że
$(Tu,v)_H=(u,g)_H$ dla wszystkich $u\in D(T)$.
Wówczas $T^*v=g$; gęstość $D(T)$ zapewnia
jednoznaczność $g$: różnica dwóch takich wektorów
jest ortogonalna do gęstego $D(T)$, a więc do całego $H$.
Jeśli $v,w\in D(T^*)$ i $a,b$ są skalarami, to
$(Tu,av+bw)_H=\overline a(u,T^*v)_H+
\overline b(u,T^*w)_H=(u,aT^*v+bT^*w)_H$
dla każdego $u\in D(T)$ (w przypadku rzeczywistym
sprzężenia są zbędne). Zatem $av+bw\in D(T^*)$ oraz
$T^*(av+bw)=aT^*v+bT^*w$.
Równoważnie $v\in D(T^*)$ wtedy i tylko wtedy, gdy
istnieje $C_v\ge0$ takie, że
$|(Tu,v)_H|\le C_v\|u\|_H$ dla każdego $u\in D(T)$.
Konieczność wynika z nierówności Cauchy'ego--Schwarza.
Dla dostateczności ograniczony funkcjonał $u\mapsto(Tu,v)_H$
przedłuża się jednoznacznie przez gęstość na $H$, a twierdzenie
Riesza o reprezentacji funkcjonału daje wektor $g\in H$
spełniający definicję sprzężenia.
Operator $T^*$ jest domknięty: jeśli $v_n\to v$
i $T^*v_n\to g$, to dla ustalonego $u\in D(T)$ przechodzimy
do granicy w $(Tu,v_n)_H=(u,T^*v_n)_H$, dostając
$(Tu,v)_H=(u,g)_H$. Z definicji $v\in D(T^*)$ i $T^*v=g$.
Operator jest \emph{symetryczny}, gdy
$(Tu,v)=(u,Tv)$ dla wszystkich $u,v\in D(T)$,
czyli $T\subset T^*$. Jest \emph{samosprzężony}, gdy
$T=T^*$, łącznie z równością dziedzin.
\end{definicja}

\begin{uwaga}
Symetria nie gwarantuje samosprzężoności; przykład
\ref{ex:pde12-derivative} pokazuje to nawet dla domkniętego
operatora różniczkowego pierwszego rzędu. Warunki brzegowe
mogą radykalnie zmienić dziedzinę operatora sprzężonego,
a więc także jego widmo i ewolucję.
\end{uwaga}

\section{Konstrukcja laplasjanu Dirichleta przez twierdzenie Laxa--Milgrama}
\label{sec:pde12-dirichlet}

W konstrukcji laplasjanu Dirichleta i jej zastosowaniach
przyjmujemy, o ile nie zaznaczono inaczej, że
$\Omega\subset\R^d$ jest niepustym obszarem ograniczonym
Lipschitzowskim, $H=L^2(\Omega;\R)$ i $V=H_0^1(\Omega;\R)$.
Przechodząc do widma zespolonego i grup unitarnych,
zespolimy te przestrzenie i operator. Przestrzeń $V$
wyposażamy w normę
$\|u\|_V^2=\|\nabla u\|_2^2+\|u\|_2^2$.
Rozważmy formę $b(u,v)=\int_\Omega(\nabla u\cdot\nabla v+uv)
\,\mathrm d\lambda_d(x)$. Jest ona iloczynem skalarnym
pełnej normy $H^1$ na $V$. Przestrzeń $V$ jest zupełna,
gdyż jest domkniętą podprzestrzenią $H^1$ na mocy
twierdzenia~\ref{thm:pde9-banach} i definicji $H_0^1$.
Jest też gęsta w $H$, ponieważ zawiera funkcje testowe.

\begin{twierdzenie}[Rezolwenta z twierdzenia Laxa--Milgrama]
\label{thm:pde12-resolvent}
Dla każdego $f\in H$ istnieje dokładnie jedno $u=Bf\in V$
takie, że $b(u,v)=(f,v)_H$ dla wszystkich $v\in V$.
Operator $B:H\to H$ jest liniowy, ograniczony, zwarty,
samosprzężony, dodatni i różnowartościowy. Ponadto
$\|B\|_{H\to H}\le1$ oraz $\overline{\operatorname{Ran}B}=H$.
\end{twierdzenie}
\begin{proof}
\textbf{Istnienie i oszacowanie.} Forma $b$ jest
iloczynem skalarnym na $V$, a więc jest ciągła i koercywna
ze stałą $1$. Funkcjonał $v\mapsto(f,v)_H$ jest ciągły na $V$:
$|(f,v)_H|\le\|f\|_2\|v\|_V$. Z twierdzenia
Laxa--Milgrama~\ref{thm:pde11-LM} otrzymujemy jedyne $u$.
Przy $v=u$ mamy
$\|u\|_V^2=(f,u)_H\le\|f\|_2\|u\|_2\le\|f\|_2\|u\|_V$.
Jeżeli $u\ne0$, dzielimy przez $\|u\|_V>0$;
gdy $u=0$, nierówność zachodzi bez dzielenia.
Stąd $\|u\|_V\le\|f\|_2$, a tym bardziej
$\|Bf\|_2\le\|f\|_2$. Dla $f,g\in H$ i $s,t\in\R$
dwuliniowość formy daje
$b(sBf+tBg,v)=(sf+tg,v)_H$ dla każdego $v\in V$.
Jednoznaczność implikuje $B(sf+tg)=sBf+tBg$.
Otrzymujemy ciągłość zarówno $B:H\to V$, jak i $B:H\to H$.

\textbf{Zwarty charakter.} Jeśli $(f_n)$ jest ograniczony w $H$,
to $(Bf_n)$ jest ograniczony w $V$. Twierdzenie
Rellicha--Kondrachova~\ref{thm:pde10-Rellich-p} (w postaci
zanurzenia $H_0^1(\Omega)\hookrightarrow L^2(\Omega)$)
dostarcza podciągu zbieżnego w $H$. To znaczy, że $B$
jest zwarty.

\textbf{Samosprzężoność i dodatniość.} Dla $u=Bf$, $w=Bg$
możemy wybrać odpowiednio $v=w$ i $v=u$ w równaniach
wariacyjnych. Symetria $b$ daje
$(Bf,g)_H=(g,u)_H=b(w,u)=b(u,w)=(f,Bg)_H$.
Ponieważ $B$ jest ograniczony i określony na całym $H$,
tożsamość ta daje $D(B^*)=H$ oraz $B^*=B$.
Dalej $(Bf,f)_H=b(Bf,Bf)=\|Bf\|_V^2\ge0$.

\textbf{Jądro i gęstość obrazu.} Gdy $Bf=0$, jego równanie
wariacyjne daje $(f,v)_H=0$ dla wszystkich $v\in V$.
Ponieważ $C_c^\infty(\Omega)\subset V$ jest gęste w $H$,
to $f=0$. Wreszcie
$y\perp\operatorname{Ran}B$ wtedy i tylko wtedy, gdy
$(Bf,y)_H=(f,By)_H=0$ dla każdego $f\in H$, czyli $By=0$.
Zatem $(\operatorname{Ran}B)^\perp=\ker B=\{0\}$.
Twierdzenie o rzucie ortogonalnym na domkniętą podprzestrzeń
(lemat~\ref{lem:pde11-projection}) daje
$\overline{\operatorname{Ran}B}=H$.
\end{proof}


\begin{uwaga}[Z czego wynika zwartość rezolwenty?]
\label{rem:pde-v42-compact-chain}
Zależność logiczna w poprzednim dowodzie jest istotna:
koercywność formy i twierdzenie Laxa--Milgrama
(\ref{thm:pde11-LM}) zapewniają oszacowanie
$B:L^2(\Omega)\to H_0^1(\Omega)$, a dopiero
zwartość zanurzenia Rellicha--Kondrachova
(\ref{thm:pde10-Rellich-p}) zapewnia zwartość
$B:L^2(\Omega)\to L^2(\Omega)$. Samo istnienie
ograniczonej rezolwenty nie daje zwartości.
W szczególności ograniczoność obszaru i warunki brzegowe
odgrywają rolę przy późniejszym rozkładzie spektralnym.
\end{uwaga}

\begin{definicja}[Laplasjan Dirichleta w $L^2$]
\label{def:pde12-laplacian}
Operatorem Dirichleta $A=-\Delta_D$ nazywamy operator w $H$
o dziedzinie
\[
D(A)=\big\{u\in V:\ \exists g\in H\quad
\int_\Omega\nabla u\cdot\nabla v\,\mathrm d\lambda_d(x)
=(g,v)_H\ \text{dla każdego }v\in V\big\},
\]
przy czym $Au=g$. Element $g$ jest wyznaczony
jednoznacznie dzięki gęstości $V$ w $H$: różnica dwóch
reprezentantów jest ortogonalna do całego $V$.
Jeśli $u_1,u_2\in D(A)$ mają reprezentantów $g_1,g_2$,
to dla skalarów $s,t\in\R$ dwuliniowość całki pokazuje,
że $su_1+tu_2\in D(A)$ z reprezentantem $sg_1+tg_2$.
Zatem $D(A)$ jest podprzestrzenią liniową i $A$ jest liniowy.
Dla każdego $u\in D(A)$, podstawiając test
$v\in C_c^\infty(\Omega)$, dostajemy
$\langle-\Delta u,v\rangle=
\int_\Omega\nabla u\cdot\nabla v\,\mathrm d\lambda_d(x)
=(Au,v)_H$. Odwrotnie, jeśli $u\in V$ i dystrybucja
$-\Delta u$ jest reprezentowana przez $g\in L^2$, to ta
sama tożsamość dla testów rozszerza się na $V$ przez
gęstość w $H^1$ i nierówność Cauchy'ego--Schwarza.
Wobec tego
$D(A)=\{u\in H_0^1(\Omega):\Delta u\in L^2(\Omega)
\text{ w sensie dystrybucji}\}$.
Warunek $u\in H_0^1(\Omega)$ narzuca zerowy ślad
(twierdzenie~\ref{thm:pde10-trace-zero}).
\end{definicja}

\begin{twierdzenie}[Samosprzężoność laplasjanu Dirichleta]
\label{thm:pde12-A-selfadjoint}
Operator $A$ z definicji~\ref{def:pde12-laplacian} jest
gęsto określony, nieujemny i samosprzężony.
Ponadto $D(A)=\operatorname{Ran}B$, a na tej dziedzinie
$A=B^{-1}-I$, więc $(A+I)^{-1}=B$.
\end{twierdzenie}
\begin{proof}
Jeśli $u=Bf$, to
$\int_\Omega\nabla u\cdot\nabla v\,\mathrm d\lambda_d(x)
=(f-u,v)_H$, czyli $u\in D(A)$ i $Au=f-u$.
Odwrotnie, jeśli $u\in D(A)$ i $Au=g$, to
$b(u,v)=(g+u,v)_H$, a z jednoznaczności równania
dla $B$ wynika $u=B(g+u)$. Tak otrzymujemy
$D(A)=\operatorname{Ran}B$, a więc gęstość dziedziny
z twierdzenia~\ref{thm:pde12-resolvent}.
W szczególności $B:H\to D(A)$ jest bijekcją,
$(A+I)Bf=f$ dla $f\in H$ i $B(A+I)u=u$ dla $u\in D(A)$.
Oznacza to $(A+I)^{-1}=B$ i $A=B^{-1}-I$ na $D(A)$.
Wybór $v=u$ w definicji daje
$(Au,u)_H=\|\nabla u\|_2^2\ge0$.

Aby wykazać samosprzężoność bez ukrywania dziedziny,
zauważmy, że $A$ jest symetryczny:
$(Au,v)_H=(u,Av)_H$ dla $u,v\in D(A)$
przez symetrię całki z gradientów. Jeśli $w\in D(A^*)$,
to dla każdego $u=Bf\in D(A)$ mamy
$(Au,w)_H=(u,A^*w)_H$, a więc
$(f-Bf,w)_H=(Bf,A^*w)_H=(f,BA^*w)_H$.
Samosprzężoność $B$ daje
$(f,w)_H=(f,Bw+BA^*w)_H$ dla wszystkich $f\in H$.
Zatem $w=B(w+A^*w)\in\operatorname{Ran}B=D(A)$.
Otrzymaliśmy $D(A^*)\subset D(A)$; inkluzja przeciwna
wynika z symetrii. Tym samym $A=A^*$. Operator jest również domknięty
jako operator sprzężony, zgodnie z definicją
~\ref{def:pde12-adjoint}. Wprost: gdy $u_n\in D(A)$,
$u_n\to u$ i $Au_n\to g$ w $H$, ciągłość $B$ daje
$u_n=B(Au_n+u_n)\to B(g+u)$, więc
$u=B(g+u)\in D(A)$ oraz $Au=g$.
\end{proof}

\begin{uwaga}[Dziedzina nie zawsze jest $H^2\cap H_0^1$]
\label{rem:pde12-corners}
Jeśli brzeg $\partial\Omega$ jest klasy $C^{1,1}$, to
$D(A)=H^2(\Omega)\cap H_0^1(\Omega)$: inkluzja z prawej
strony wynika z dystrybucyjnej charakterystyki $D(A)$,
a odwrotna z globalnej regularności eliptycznej
(twierdzenie~\ref{thm:pde15-boundary} dla macierzy jednostkowej).
Nie korzystamy z regularności $H^2$ w dowodzie samosprzężoności.

Dla ogólnego obszaru Lipschitzowskiego równość może zawieść.
Niech $\Omega=\{(r\cos\theta,r\sin\theta):0<r<1,\
0<\theta<\omega\}$, gdzie $\pi<\omega<2\pi$, i
$\beta=\pi/\omega\in(1/2,1)$. Jest to obszar Lipschitzowski:
przy wierzchołku wybierzmy oś $y$ wzdłuż dwusiecznej
dopełniającego sektora o kącie $\gamma=2\pi-\omega\in(0,\pi)$.
W tych współrzędnych, w dostatecznie małym otoczeniu zera,
obszar leży poniżej wykresu $y=|x|\cot(\gamma/2)$,
który jest Lipschitzowski. Poza wierzchołkiem brzeg
składa się z gładkich łuków i odcinków, a w punktach
ich przecięcia kąt wewnętrzny wynosi $\pi/2$, więc
po obrocie także tam otrzymujemy lokalny wykres
Lipschitzowski.
Wybierzmy gładkie odcięcie radialne $\eta$, równe $1$
dla $0\le r\le1/3$ i $0$ dla $r\ge2/3$, oraz połóżmy
$w(r,\theta)=r^\beta\sin(\beta\theta)$ i $u=\eta w$.
Dla $r>0$ mamy
\[
 \begin{aligned}
 \Delta w&=\partial_r^2 w+r^{-1}\partial_r w+
 r^{-2}\partial_\theta^2 w\\
 &= [\beta(\beta-1)+\beta-\beta^2]r^{\beta-2}
 \sin(\beta\theta)=0,\\
 |\nabla w|^2&=|\partial_r w|^2+r^{-2}|\partial_\theta w|^2
 =\beta^2r^{2\beta-2}.
 \end{aligned}
\]
W pobliżu wierzchołka zarówno $u$, jak i $\nabla u$
należą do $L^2$, ponieważ w biegunowych współrzędnych
ich kwadraty całkują się z czynnikiem $r$ i wykładniki
$2\beta+1$, $2\beta-1$ są większe od $-1$.
Z dala od wierzchołka funkcja jest gładka z ograniczonymi
pochodnymi. Pochodne klasyczne we wnętrzu są słabymi
pochodnymi, więc $u\in H^1(\Omega)$.

Aby uzasadnić zerowy ślad, wybierzmy radialne gładkie
odcięcie $\eta_\varepsilon$ znikające dla $r\le\varepsilon$,
równe $1$ dla $r\ge2\varepsilon$ i spełniające
$|\eta_\varepsilon'|\le C/\varepsilon$ (przez przeskalowanie
ustalonej funkcji). Funkcje $u_\varepsilon=\eta_\varepsilon u$
są gładkie aż do brzegu i znikają na nim, więc mają
zerowy ślad. Zbieżność $u_\varepsilon\to u$ w $H^1$
wynika z całkowalności $u,\nabla u$ oraz oszacowania
\[
 \|u\nabla\eta_\varepsilon\|_2^2
 \le C_1\varepsilon^{-2}\int_\varepsilon^{2\varepsilon}
 r^{2\beta+1}\,\mathrm d\lambda_1(r)
 \le C_2\varepsilon^{2\beta}\longrightarrow0.
\]
Z ciągłości śladu i twierdzenia~\ref{thm:pde10-trace-zero}
otrzymujemy $u\in H_0^1(\Omega)$.
Ponieważ $\Delta u=2\nabla\eta\cdot\nabla w+w\Delta\eta$
ma nośnik z dala od wierzchołka i jest ograniczona,
$\Delta u\in L^2(\Omega)$. Nie występuje dodatkowy składnik
dystrybucyjny w wierzchołku, który nie należy do $\Omega$.
Z definicji~\ref{def:pde12-laplacian} mamy $u\in D(A)$.

Tymczasem dla $r<1/3$ zachodzi
$\partial_r^2u=\beta(\beta-1)r^{\beta-2}\sin(\beta\theta)$.
Na podsektorze $\omega/4<\theta<3\omega/4$ moduł sinusa
jest nie mniejszy niż $1/\sqrt2$. Ponieważ dla jednostkowego
wektora radialnego $e_r$ mamy
$\partial_r^2u=e_r^T(\nabla^2u)e_r$,
moduł tej pochodnej nie przekracza normy Frobeniusa
macierzy drugich pochodnych. Wobec tego
\[
 \int_{\Omega\cap B(0,1/3)}|\nabla^2u|^2\,\mathrm d\lambda_2
 \ge c\int_0^{1/3}r^{2\beta-3}\,\mathrm d\lambda_1(r)
 =+\infty,
\]
gdyż $\beta<1$ i $c>0$. Gdyby $u\in H^2$, jej słabe
drugie pochodne zgadzałyby się we wnętrzu z klasycznymi,
wbrew tej rozbieżności. Zatem $u\notin H^2(\Omega)$.
Szczególny przypadek $\omega=3\pi/2$ powróci w
kontrprzykładzie~\ref{ex:pde15-corner}.
\end{uwaga}

\section{Twierdzenie spektralne dla zwartego operatora dodatniego}
\label{sec:pde12-compact-spectrum}

\begin{twierdzenie}[Rozkład spektralny zwartego operatora dodatniego]
\label{thm:pde12-compact-spectral}
Niech $H$ będzie rzeczywistą lub zespoloną,
separowalną przestrzenią Hilberta, a $K:H\to H$
operatorem zwartym, samosprzężonym i nieujemnym.
Wówczas istnieje ortonormalny układ wektorów własnych
$(e_j)$ odpowiadających dodatnim wartościom własnym
$\mu_j$, uporządkowanym nierosnąco i z uwzględnieniem
krotności. Układ może być pusty lub skończony.
Każda dodatnia wartość własna ma skończoną krotność.
Jeśli układ jest nieskończony, to $\mu_j\to0$. Ponadto
$H=\ker K\oplus\overline{\operatorname{span}}\{e_j\}$,
a dla $x\in H$ zachodzi
$Kx=\sum_j\mu_j(x,e_j)e_j$ zbieżnie w $H$.
Jeżeli $K$ jest różnowartościowy, $(e_j)$ jest bazą
ortonormalną całego $H$.
\end{twierdzenie}
\begin{proof}
Jeżeli $K=0$ (w szczególności gdy $H=\{0\}$), wybieramy
pusty układ, a wszystkie twierdzenia wynikają z $\ker K=H$.
Zakładamy dalej $K\ne0$.

\textbf{Krok 1: maksimum ilorazu Rayleigha.}
Połóżmy $m=\sup_{\|x\|=1}(Kx,x)$.
Pokażemy, że $m=\|K\|>0$. Z nierówności
Cauchy'ego--Schwarza w $H$ mamy
$0\le (Kx,x)\le\|Kx\|\|x\|\le\|K\|$ dla $\|x\|=1$,
więc $0\le m\le\|K\|$.
Dodatniość daje nierówność Cauchy'ego--Schwarza
dla formy $(Kx,y)$:
$|(Kx,y)|^2\le(Kx,x)(Ky,y)$.
Istotnie, dla $q(x,y)=(Kx,y)$ forma $q$ jest hermitowska
(w przypadku rzeczywistym symetryczna).
Niech $a=q(x,y)$ i $b=q(y,y)\ge0$.
Dla dowolnego skalara $t$ mamy
$q(x-ty,x-ty)=q(x,x)-2\operatorname{Re}(\overline t a)
+|t|^2b\ge0$. Gdy $b>0$, wybór $t=a/b$ daje żądaną
nierówność; gdy $b=0$, wybór $t=sa$ dla dowolnego
$s>0$ wymusza $a=0$.
Dla wektorów jednostkowych $x,y$ otrzymujemy
$|(Kx,y)|\le m$; dla $Kx\ne0$ wybieramy
$y=Kx/\|Kx\|$, otrzymując $\|Kx\|\le m$.
Dla $Kx=0$ ta nierówność również zachodzi.
Biorąc supremum po jednostkowych $x$, dostajemy
$\|K\|\le m$. Zatem $m=\|K\|>0$, ponieważ $K\ne0$.
Wybierzmy $\|x_n\|=1$ z $(Kx_n,x_n)\to m$.
Z ograniczoności i słabej zwartości ciągowej kul
(lemat~\ref{lem:pde11-weakcompact} w przypadku rzeczywistym;
dla zespolonego stosujemy go do przestrzeni rzeczywistej
z iloczynem $\operatorname{Re}(\cdot,\cdot)$: testowanie
wektorami $y$ i $iy$ daje zbieżność odpowiednio części
rzeczywistych i urojonych iloczynów skalarnych)
wybieramy podciąg $x_n\rightharpoonup x$,
zachowując oznaczenie indeksów.
Mamy $\|x\|\le1$, ponieważ dla $x\ne0$
$\|x\|^2=\lim_n |(x_n,x)|\le\|x\|$,
a dla $x=0$ nierówność jest spełniona.
Zwarty operator przenosi ten ciąg w ciąg silnie zbieżny:
gdyby $Kx_n\not\to Kx$, pewien podciąg miałby odległość
od $Kx$ co najmniej $\varepsilon>0$. Z jego zwartości
wybieramy dalszy podciąg zbieżny w normie do $z$.
Dla każdego $y\in H$ zachodzi
$(z,y)=\lim_k(x_{n_k},Ky)=(x,Ky)=(Kx,y)$,
więc $z=Kx$, sprzeczność. Ponadto
\[
 |(Kx_n,x_n)-(Kx,x)|\le
 \|Kx_n-Kx\|\|x_n\|+|(Kx,x_n-x)|\longrightarrow0.
\]
Stąd $(Kx,x)=m$.
Ponieważ $(Kx,x)\le m\|x\|^2$ i $\|x\|\le1$,
otrzymujemy $\|x\|=1$.

\textbf{Krok 2: maksimum daje wektor własny.}
Dla $y\perp x$ rozpatrujemy
$\varphi(t)=\frac{(K(x+ty),x+ty)}{\|x+ty\|^2}$,
gdzie $t\in\R$. Ortogonalność daje mianownik
$1+t^2\|y\|^2$, zaś samosprzężoność daje licznik
$m+2t\operatorname{Re}(Kx,y)+t^2(Ky,y)$.
Funkcja osiąga maksimum przy $t=0$, więc
$0=\varphi'(0)=2\operatorname{Re}(Kx,y)$.
W przypadku zespolonym powtarzamy argument dla $iy$:
$\operatorname{Re}(Kx,iy)=\operatorname{Re}(-i(Kx,y))
=\operatorname{Im}(Kx,y)$, więc także część urojona
jest zerowa. Zatem $Kx-mx$ jest ortogonalny do
$x^\perp$ oraz do $x$, ponieważ $(Kx,x)=m$ i $\|x\|=1$.
W konsekwencji $Kx-mx=0$.

\textbf{Krok 3: iteracja na ortogonalnym dopełnieniu.}
Przestrzeń $x^\perp$ jest niezmiennicza względem $K$,
bo $(Ky,x)=(y,Kx)=m(y,x)=0$. Ograniczenie $K$ do niej
pozostaje zwarte, samosprzężone i dodatnie. Powtarzamy
konstrukcję, dopóki ograniczenie nie stanie się zerowe.
Ściślej, po wyborze $e_1,\ldots,e_j$ rozpatrujemy
$M_j=\operatorname{span}\{e_1,\ldots,e_j\}^\perp$.
Jest to domknięta podprzestrzeń niezmiennicza, gdyż
dla $y\in M_j$ mamy $(Ky,e_\ell)=\mu_\ell(y,e_\ell)=0$
dla $\ell=1,\ldots,j$. Ograniczenie $K|_{M_j}$
jest zwarte (granice podciągów leżą w domkniętym $M_j$),
nieujemne i samosprzężone jako operator na całym $M_j$.
Jeśli $K|_{M_j}\ne0$, stosujemy kroki 1--2 w $M_j$
i otrzymujemy $\mu_{j+1}=\|K|_{M_j}\|>0$ i jednostkowy
$e_{j+1}\in M_j$. Ponieważ $M_j\subset M_{j-1}$,
wartości te są nierosnące, a wektory ortonormalne.
W wymiarze skończonym konstrukcja kończy się najpóźniej
po $\dim H$ krokach.
Jeśli układ jest nieskończony i $\mu_j$ nie dążą do zera,
istnieje $c>0$ i nieskończenie wiele indeksów, dla których
$\|Ke_j\|=\mu_j\ge c$. Dla dwóch różnych takich indeksów
$\|Ke_j-Ke_k\|^2=\mu_j^2+\mu_k^2\ge2c^2$.
Zatem obrazy odpowiedniego podciągu jednostkowych $e_j$
nie mają podciągu Cauchy'ego, co przeczy zwartości $K$.

\textbf{Krok 4: zupełność na $(\ker K)^\perp$.}
Jeżeli $z\perp e_j$ dla wszystkich $j$, to $z$
należy do kolejnych ortogonalnych dopełnień, na których
norma ograniczenia $K$ jest co najwyżej $\mu_{j+1}$.
Zatem $\|Kz\|\le\mu_{j+1}\|z\|\to0$;
gdy konstrukcja skończyła się po skończonej liczbie kroków,
operator na ostatnim dopełnieniu jest z definicji zerowy.
W obu przypadkach $Kz=0$. Odwrotnie, jeśli $Kz=0$,
to $\mu_j(z,e_j)=(z,Ke_j)=(Kz,e_j)=0$, więc $z$
jest ortogonalny do wszystkich $e_j$. Otrzymujemy
$\ker K=(\overline{\operatorname{span}}\{e_j\})^\perp$,
co daje rozkład ortogonalny. Dla $x\in H$ rozwinięcie
$x=x_0+\sum_j(x,e_j)e_j$ z $x_0\in\ker K$ jest
zbieżne w $H$; stosując ograniczony operator $K$
do sum częściowych, uzyskujemy żądany wzór.
Jeśli $Ky=\mu y$ dla $\mu>0$, zapiszmy
$y=y_0+\sum_j(y,e_j)e_j$, gdzie $y_0\in\ker K$.
Rzut równania na $\ker K$ daje $0=\mu y_0$, więc $y_0=0$.
Sparowanie równania z $e_j$ daje
$(\mu_j-\mu)(y,e_j)=0$ dla wszystkich $j$.
Zatem przestrzeń własna jest rozpięta przez dokładnie
te $e_j$, dla których $\mu_j=\mu$; jest skończeniewymiarowa,
gdyż $\mu_j\to0$ w przypadku nieskończonego układu.
Przy $\ker K=\{0\}$ układ jest bazą całego $H$.
\end{proof}

\begin{twierdzenie}[Twierdzenie spektralne dla zwartego operatora samosprzężonego]
\label{thm:pde12-compact-general-spectral}
Niech $H$ będzie rzeczywistą lub zespoloną separowalną
przestrzenią Hilberta, a $K\in\mathcal L(H)$ operatorem
zwartym i samosprzężonym, niekoniecznie dodatnim.
Wówczas $H$ jest domkniętą sumą ortogonalną $\ker K$ oraz
podprzestrzeni własnych odpowiadających niezerowym
rzeczywistym wartościom własnym $\kappa_j$.
Każda niezerowa wartość własna ma skończoną krotność,
a jeśli jest ich nieskończenie wiele (licząc z krotnościami),
to $|\kappa_j|\to0$. Istnieje ortonormalny układ
$(e_j)$ tych wektorów własnych taki, że
\[
 Kx=\sum_j\kappa_j(x,e_j)_He_j\quad(x\in H),
\]
gdzie szereg jest zbieżny w $H$; jeżeli $\ker K=\{0\}$,
to $(e_j)$ jest bazą ortonormalną całego $H$.
\end{twierdzenie}
\begin{proof}
Operator $K^2$ jest zwarty, samosprzężony i nieujemny:
$K^2$ jest zwarty jako złożenie operatora zwartego
i ograniczonego, a $(K^2x,x)=\|Kx\|^2\ge0$.
Ponadto $\ker K^2=\ker K$: jeśli $K^2x=0$, to
$\|Kx\|^2=(K^2x,x)=0$; druga inkluzja jest bezpośrednia.
Twierdzenie~\ref{thm:pde12-compact-spectral} zastosowane
do $K^2$ daje rozkład ortogonalny
$H=\ker K\oplus\overline{\bigoplus_{\mu>0}E_\mu}$,
gdzie $E_\mu=\ker(K^2-\mu I)$, $\mu>0$.
Każde $E_\mu$ jest skończeniewymiarowe, gdyż nieskończony
układ ortonormalny w $E_\mu$ miałby obrazy przez $K^2$
bez podciągu zbieżnego, co przeczyłoby zwartości.

Ponieważ $K^2(Kx)=K(K^2x)=\mu Kx$ dla $x\in E_\mu$,
każda przestrzeń $E_\mu$ jest niezmiennicza względem $K$.
Na $E_\mu$ operator $J=\mu^{-1/2}K|_{E_\mu}$
jest samosprzężony i spełnia $J^2=I_{E_\mu}$.
Każdy $x\in E_\mu$ ma rozkład
$x=x_++x_-$, gdzie $x_\pm=(x\pm Jx)/2$.
Bezpośrednio $Jx_\pm=\pm x_\pm$, więc
$Kx_\pm=\pm\sqrt\mu\,x_\pm$.
Podprzestrzenie odpowiadające obu znakom są ortogonalne:
dla $Jv=v$ i $Jw=-w$ mamy
$(v,w)=(Jv,w)=(v,Jw)=-(v,w)$.
W każdej z nich wybieramy bazę ortonormalną
(może być pusta). W ten sposób otrzymujemy bazę własną
całego $E_\mu$ bez odwołania do odrębnego twierdzenia
o diagonalizacji macierzy.
Bazy wszystkich $E_\mu$ tworzą wraz z bazą jądra
ortogonalny rozkład $H$ na wektory własne $K$.
Każda niezerowa wartość własna $\kappa$ operatora $K$
ma przestrzeń własną zawartą w skończeniewymiarowym
$E_{\kappa^2}$. Dla każdego $\varepsilon>0$ istnieje
tylko skończenie wiele wektorów w uzyskanej bazie własnej
z $|\kappa_j|\ge\varepsilon$, ponieważ odpowiadają
wartościom własnym $K^2$ nie mniejszym niż
$\varepsilon^2$, liczonym z krotnościami.
Jeśli układ niezerowych modów jest nieskończony,
to $|\kappa_j|\to0$.
Dla sum skończonych podany wzór na $K$ wynika
z diagonalizacji. Dla dowolnego $x\in H$ przechodzimy
do granicy, ponieważ rzuty na skończone sumy podprzestrzeni
własnych zbiegają silnie do rzutu na $(\ker K)^\perp$,
a $K$ jest ograniczony i znika na swoim jądrze.
\end{proof}

\begin{wniosek}[Dyskretne widmo laplasjanu Dirichleta]
\label{cor:pde12-eigen}
Dla obszaru ograniczonego Lipschitzowskiego $\Omega$ operator
$A=-\Delta_D$ ma bazę ortonormalną $(\phi_j)_{j\ge1}$ w
$L^2(\Omega)$ złożoną z wektorów własnych:
$A\phi_j=\lambda_j\phi_j$, gdzie
$0<\lambda_1\le\lambda_2\le\cdots$, $\lambda_j\to\infty$.
Każda wartość własna ma skończoną krotność,
a w sensie słabym
$\int_\Omega\nabla\phi_j\cdot\nabla v\,\mathrm d\lambda_d(x)
=\lambda_j\int_\Omega\phi_jv\,\mathrm d\lambda_d(x)$
dla wszystkich $v\in H_0^1(\Omega)$.
\end{wniosek}
\begin{proof}
Dla operatora $B=(A+I)^{-1}$ z
\ref{thm:pde12-resolvent} jądro jest zerowe,
więc twierdzenie~\ref{thm:pde12-compact-spectral}
daje bazę $(\phi_j)$ i liczby $\mu_j>0$ z
$B\phi_j=\mu_j\phi_j$.
Przestrzeń $L^2(\Omega)$ jest separowalna i
nieskończeniewymiarowa: w kuli zawartej w niepustym
$\Omega$ można wybrać nieskończenie wiele parami
rozłącznych mniejszych kul i znormalizować ich funkcje
charakterystyczne, uzyskując układ ortonormalny.
Ponieważ $B$ jest różnowartościowy, baza nie może
być skończona, a $\mu_j\downarrow0$ (licząc krotności).
Ponadto $\phi_j=\mu_j^{-1}B\phi_j\in\operatorname{Ran}B=D(A)$.
Z równości $A=B^{-1}-I$ otrzymujemy
$A\phi_j=(\mu_j^{-1}-1)\phi_j$, czyli
$\lambda_j=\mu_j^{-1}-1$. Nieujemność $A$ daje
$\lambda_j\ge0$. Gdyby $\lambda_j=0$, to
$\|\nabla\phi_j\|_2^2=(A\phi_j,\phi_j)=0$.
Nierówność Poincarégo w $H_0^1(\Omega)$
(wniosek~\ref{cor:pde10-Poincare} dla $p=2$) daje
$\phi_j=0$, sprzeczność z normalizacją.
Zatem $\lambda_j>0$, a $\lambda_j\to+\infty$,
ponieważ $\mu_j\to0$. Dla $\lambda\ge0$ równanie
$Au=\lambda u$ z $u\ne0$ jest równoważne
$Bu=(1+\lambda)^{-1}u$: w drugim kierunku
$u=(1+\lambda)Bu\in D(A)$. Zatem krotności są równe
krotnościom odpowiednich dodatnich wartości własnych
$B$, skończonym na mocy twierdzenia spektralnego.
Tożsamość słaba wynika z definicji $A$.
Brak innych punktów widma $A$ niż wartości własne
wykażemy w twierdzeniu~\ref{thm:pde12-fredholm}.
\end{proof}

\section{Dziedzina formy, zasada min--max i oszacowania widma}
\label{sec:pde12-minmax}

\begin{twierdzenie}[Rozwinięcie spektralne formy energii]
\label{thm:pde12-formdomain}
Niech $(\lambda_j,\phi_j)$ będą wartościami i funkcjami
własnymi z wniosku~\ref{cor:pde12-eigen}. Dla $u\in L^2(\Omega)$
niech $c_j=(u,\phi_j)_{L^2}$. Wówczas
\begin{align*}
u\in H_0^1(\Omega)
&\quad\Longleftrightarrow\quad
\sum_{j=1}^\infty\lambda_j|c_j|^2<\infty,\\
u\in D(A)
&\quad\Longleftrightarrow\quad
\sum_{j=1}^\infty\lambda_j^2|c_j|^2<\infty.
\end{align*}
Dla $u\in H_0^1(\Omega)$ zachodzi
$\|\nabla u\|_2^2=\sum_j\lambda_j|c_j|^2$,
a dla $u\in D(A)$ zachodzi
$Au=\sum_j\lambda_jc_j\phi_j$ zbieżnie w $L^2$.
\end{twierdzenie}
\begin{proof}
Piszemy $\mu_j=(1+\lambda_j)^{-1}$.
Ponieważ $\phi_j=\mu_j^{-1}B\phi_j$, mamy
$\phi_j\in H_0^1(\Omega)$. Równanie wariacyjne dla $B$
daje $b(\phi_i,v)=\mu_i^{-1}(\phi_i,v)_{L^2}$
dla każdego $v\in V$. W szczególności
$b(\sqrt{\mu_i}\phi_i,\sqrt{\mu_j}\phi_j)
=\sqrt{\mu_j/\mu_i}\,\delta_{ij}=\delta_{ij}$.
Wektory $\sqrt{\mu_j}\phi_j$ tworzą więc układ
ortonormalny w $V=H_0^1(\Omega)$ z iloczynem skalarnym $b$. Układ ten jest zupełny w $V$:
jeśli $v\in V$ jest do niego ortogonalny, to
$0=b(v,\phi_j)=\mu_j^{-1}(v,\phi_j)_{L^2}$
dla wszystkich $j$; zupełność $(\phi_j)$ w $L^2$ daje $v=0$.
Z twierdzenia Parsevala w $V$ dostajemy
\[
\|v\|_V^2=\sum_j|b(v,\sqrt{\mu_j}\phi_j)|^2
=\sum_j\mu_j^{-1}|(v,\phi_j)_{L^2}|^2
=\sum_j(1+\lambda_j)|(v,\phi_j)_{L^2}|^2.
\]
Dla ogólnego $u\in L^2$ z Parsevala w $L^2$ mamy
$\sum_j|c_j|^2=\|u\|_2^2<\infty$, więc warunek
$\sum_j\lambda_j|c_j|^2<\infty$ jest równoważny
sumowalności z wagami $1+\lambda_j$.
Dla $u_N=\sum_{j\le N}c_j\phi_j$ i $M>N$ ortogonalność
w $V$ daje
$\|u_M-u_N\|_V^2=\sum_{j=N+1}^M(1+\lambda_j)|c_j|^2$.
Stąd sumy częściowe są ciągiem Cauchy'ego w $V$. Ich granica $v\in V$
równa się $u$, gdyż zbieżność w $V$ implikuje zbieżność
w $L^2$, a w $L^2$ mamy $u_N\to u$ z Parsevala.
Tak wykazaliśmy równoważność dla $H_0^1$.
Dla $u\in V$ stosujemy powyższą tożsamość Parsevala
z $v=u$ i odejmujemy $\|u\|_2^2=\sum_j|c_j|^2$,
otrzymując równość energii.
Ponadto dla $u,v\in V$, z tożsamości polaryzacyjnej
lub rozwinięcia w bazie $V$, wynika
\[
 \int_\Omega\nabla u\cdot\nabla v\,\mathrm d\lambda_d
 =\sum_j\lambda_j(u,\phi_j)_H(v,\phi_j)_H .
\]
Szereg jest bezwzględnie zbieżny z nierówności
Cauchy'ego--Schwarza dla ważonych współczynników.
W zespoleniu drugie czynniki po prawej stronie oraz
gradient $v$ po lewej zastępujemy ich sprzężeniami.

Jeśli $u\in D(A)$, to $Au\in L^2$, a dla każdego $j$
$(Au,\phi_j)=(u,A\phi_j)=\lambda_jc_j$.
Parseval daje sumowalność $\lambda_j^2|c_j|^2$ i wzór
na $Au$. Odwrotnie, jeśli ta suma jest skończona,
to $\sum_j\lambda_j|c_j|^2
\le\lambda_1^{-1}\sum_j\lambda_j^2|c_j|^2<\infty$
dzięki $\lambda_1>0$, więc $u\in H_0^1$ z pierwszej
części twierdzenia. Niech $g=\sum_j\lambda_jc_j\phi_j\in L^2$.
Dla $v=\phi_j$ symetria formy i równanie własne
z testem $u$ dają
$\int_\Omega\nabla u\cdot\nabla\phi_j\,\mathrm d\lambda_d
=\lambda_j(u,\phi_j)_{L^2}=(g,\phi_j)_{L^2}$.
Z liniowości równość zachodzi dla skończonych sum.
Takie sumy są gęste w $V$, ponieważ
$\sqrt{\mu_j}\phi_j$ tworzą bazę ortonormalną $V$.
Jeżeli $v_n\to v$ w $V$, to błędy obu stron nie przekraczają
$\|\nabla u\|_2\|\nabla(v_n-v)\|_2$ i
$\|g\|_2\|v_n-v\|_2$. Obie liczby dążą do zera.
Dostajemy równość dla każdego $v\in V$, więc
$u\in D(A)$ i $Au=g$.
\end{proof}

\begin{twierdzenie}[Zasada Couranta--Fischera]
\label{thm:pde12-minmax}
Przy założeniach wniosku~\ref{cor:pde12-eigen},
dla każdej liczby całkowitej $k\ge1$
\[
\lambda_k=
\min_{\substack{E\subset H_0^1(\Omega)\\\dim E=k}}
\ \max_{\substack{u\in E\\u\ne0}}
\frac{\displaystyle\int_\Omega|\nabla u|^2\,\mathrm d\lambda_d(x)}
{\displaystyle\int_\Omega|u|^2\,\mathrm d\lambda_d(x)}.
\]
Minimum jest osiągane dla $E=\operatorname{span}
\{\phi_1,\ldots,\phi_k\}$.
\end{twierdzenie}
\begin{proof}
Dla dowolnej $k$-wymiarowej podprzestrzeni $E\subset V$
iloraz Rayleigha można ograniczyć do sfery
$\{u\in E:\|u\|_2=1\}$. Jest ona zwarta w $E$,
ponieważ wszystkie normy w skończonym wymiarze są
równoważne. Funkcja $u\mapsto\|\nabla u\|_2^2$ jest
ciągła w normie $V$, więc maksimum na tej sferze
jest osiągane.
Dla wskazanego $E$ i $u=\sum_{j\le k}c_j\phi_j$
twierdzenie~\ref{thm:pde12-formdomain} daje
$\|\nabla u\|_2^2=\sum_{j\le k}\lambda_j|c_j|^2
\le\lambda_k\sum_{j\le k}|c_j|^2=\lambda_k\|u\|_2^2$.
Wektor $u=\phi_k$ daje równość, więc maksimum na tym $E$
wynosi $\lambda_k$. Z drugiej strony weźmy dowolne
$k$-wymiarowe $E\subset H_0^1$. Odwzorowanie liniowe
$E\to\R^{k-1}$, $u\mapsto((u,\phi_j)_H)_{j<k}$,
ma jądro wymiaru co najmniej $1$ na mocy twierdzenia
o wymiarze jądra i obrazu (dla $k=1$ przestrzeń docelowa
jest zerowa). Wybieramy $u\ne0$ z tego jądra.
Dla takiego $u$ w rozwinięciu $L^2$ znikają wszystkie
współczynniki o indeksach mniejszych od $k$, wobec czego
$\|\nabla u\|_2^2=\sum_{j\ge k}\lambda_j|c_j|^2
\ge\lambda_k\sum_{j\ge k}|c_j|^2=\lambda_k\|u\|_2^2$.
Maksimum na dowolnym $E$ jest zatem co najmniej $\lambda_k$.
Łącząc oba oszacowania, otrzymujemy wzór.
\end{proof}

\begin{wniosek}[Pierwsza wartość własna, Poincaré i monotoniczność]
\label{cor:pde12-first}
Jeśli $\Omega\subset\R^d$ jest obszarem ograniczonym Lipschitzowskim, to
$\lambda_1=\inf_{u\in H_0^1(\Omega)\setminus\{0\}}
\|\nabla u\|_2^2/\|u\|_2^2$.
Najlepsza stała $C_P$ w nierówności
$\|u\|_2\le C_P\|\nabla u\|_2$ wynosi
$C_P=\lambda_1^{-1/2}$. Ponadto jeżeli
$\Omega_1\subset\Omega_2$ są ograniczonymi obszarami
Lipschitzowskimi, to
$\lambda_k(\Omega_2)\le\lambda_k(\Omega_1)$ dla każdego $k$.
\end{wniosek}
\begin{proof}
Pierwsza równość wynika z twierdzenia
~\ref{thm:pde12-minmax} dla $k=1$, gdyż iloraz jest
stały na każdej jednowymiarowej podprzestrzeni
poza zerem. Daje ona $\|u\|_2\le
\lambda_1^{-1/2}\|\nabla u\|_2$ dla $u\ne0$,
a dla $u=0$ nierówność też zachodzi.
Wektor $\phi_1$ osiąga równość, więc żadna mniejsza
stała nie może wystarczyć.
Dla monotoniczności rozważmy przedłużenie przez zero:
funkcja z $C_c^\infty(\Omega_1)$ po przedłużeniu jest
w $C_c^\infty(\Omega_2)$ i zachowuje normy $L^2$
oraz gradientu. Dla $u\in H_0^1(\Omega_1)$ wybieramy
testy $u_n\to u$ w $H^1(\Omega_1)$. Ich przedłużenia
tworzą ciąg Cauchy'ego w $H_0^1(\Omega_2)$, a granica
w $L^2$ jest przedłużeniem $u$ przez zero. Przechodząc
do granicy w obu tożsamościach norm, otrzymujemy
liniowe izometryczne zanurzenie
$H_0^1(\Omega_1)\hookrightarrow H_0^1(\Omega_2)$,
zachowujące także osobno normę gradientu.
Każda $k$-wymiarowa podprzestrzeń pierwszej przestrzeni
staje się więc dopuszczalną podprzestrzenią drugiej,
z tym samym maksimum ilorazu Rayleigha. Przyjmując minima
otrzymujemy podaną nierówność.
\end{proof}

\begin{przyklad}[Odcinek i prostokąt]
Na odcinku $(0,L)$ rozwiązania problemu
$-u''=\lambda u$, $u(0)=u(L)=0$ mają postać
$\phi_n(x)=\sqrt{2/L}\sin(n\pi x/L)$ oraz
$\lambda_n=(n\pi/L)^2$, $n\ge1$.
Wynik dla klasycznego zagadnienia brzegowego wraz
z zupełnością bazy sinusowej udowodniono
w twierdzeniu~\ref{thm:pde6-spectrum}.
Funkcje $\phi_n$ należą do $D(A)$ na mocy charakterystyki
dystrybucyjnej i zerowego śladu. Ponieważ tworzą pełną
bazę $L^2(0,L)$, nie ma dodatkowych wartości własnych:
sparowanie równania $Au=\lambda u$ z $\phi_n$ daje
$(\lambda_n-\lambda)(u,\phi_n)_H=0$ dla każdego $n$.
Dla $u\ne0$ co najmniej jeden współczynnik jest niezerowy.
W prostokącie $(0,L_1)\times(0,L_2)$ funkcje
\[
 \phi_{m,n}(x,y)=\frac{2}{\sqrt{L_1L_2}}
 \sin(m\pi x/L_1)\sin(n\pi y/L_2),\qquad m,n\ge1,
\]
mają zerowy ślad i należą do $D(A)$, ponieważ są
gładkie na domknięciu prostokąta i ich laplasjany są
w $L^2$. Dwukrotne różniczkowanie w obu zmiennych daje
$\lambda_{m,n}=\pi^2(m^2/L_1^2+n^2/L_2^2)$.
Ortogonalność i normalizacja wynikają z twierdzenia
Fubiniego i jednowymiarowych tożsamości
$\int_0^{L_i}\sin(m\pi x/L_i)\sin(n\pi x/L_i)
\,\mathrm d\lambda_1(x)=(L_i/2)\delta_{mn}$.
Stąd iloczyn skalarny dwóch znormalizowanych funkcji
prostokąta wynosi $\delta_{mm'}\delta_{nn'}$. Aby sprawdzić zupełność,
niech $f\in L^2$ będzie ortogonalna do wszystkich
$\phi_{m,n}$. Dla każdego $m$ funkcja
$g_m(y)=\int_0^{L_1} f(x,y)\sin(m\pi x/L_1)\,
\mathrm d\lambda_1(x)$ należy do $L^2(0,L_2)$
z nierówności Cauchy'ego--Schwarza i twierdzenia Fubiniego.
Z ortogonalności $f$ wynika, że wszystkie współczynniki
sinusowe $g_m$ są zerowe, zatem $g_m=0$ prawie wszędzie.
Po odrzuceniu sumy przeliczalnie wielu zbiorów miary zero
dla prawie każdego $y$ wszystkie współczynniki funkcji
$f(\cdot,y)\in L^2(0,L_1)$ znikają. Zupełność bazy
jednowymiarowej daje $f(\cdot,y)=0$ prawie wszędzie,
a Fubini daje $f=0$ w $L^2$ prostokąta.
W konsekwencji wymienione funkcje tworzą pełną bazę własną.
Wartości te należy uporządkować rosnąco z krotnościami.
Na kwadracie o boku $L$ pierwsza wartość wynosi
$2\pi^2/L^2$ i jest prosta, a druga i trzecia są równe
$5\pi^2/L^2$ (pary $(1,2)$ i $(2,1)$).
Dalsze krotności mogą wynikać nie tylko z symetrii
$\lambda_{m,n}=\lambda_{n,m}$, lecz także z różnych
przedstawień liczby jako sumy dwóch dodatnich kwadratów.
\end{przyklad}

\section{Widmo, rezolwenta i alternatywa Fredholma}
\label{sec:pde12-fredholm}

\begin{definicja}[Zbiór rezolwenty i widmo]
\label{def:pde12-spectrum}
Dla domkniętego gęsto określonego operatora $T$ w
zespolonej przestrzeni Hilberta jego \emph{zbiorem rezolwenty}
$\rho(T)$ nazywamy zbiór liczb $z\in\C$, dla których
$T-zI:D(T)\to H$ jest bijekcją, a odwrotność
$(T-zI)^{-1}:H\to H$ jest ograniczona.
\emph{Widmo} to $\sigma(T)=\C\setminus\rho(T)$.
Liczba $z$ jest wartością własną, jeżeli
$0\ne u\in D(T)$ spełnia $Tu=zu$.
Jeżeli $T-zI$ jest bijekcją, ograniczoność odwrotności
wynika już z twierdzenia o wykresie domkniętym:
z $f_n\to f$ i $u_n=(T-zI)^{-1}f_n\to u$ wynika
$Tu_n=f_n+zu_n\to f+zu$. Domkniętość $T$ daje
$u\in D(T)$ i $(T-zI)u=f$, czyli $u=(T-zI)^{-1}f$.
Wykres odwrotności jest zatem domknięty w $H\times H$.
Jest ona liniowa i określona na całej przestrzeni Banacha
$H$, więc twierdzenie o wykresie domkniętym daje
jej ograniczoność. Ponadto dla $R_z=(T-zI)^{-1}$ mamy
$TR_zf=f+zR_zf$, skąd
\[
 \|R_zf\|_{D(T)}^2
 \le\bigl(\|R_z\|^2+(1+|z|\|R_z\|)^2\bigr)\|f\|^2.
\]
Zatem $R_z:H\to D(T)$ jest ograniczony również
w normie grafowej. Wyróżniamy ponadto trzy rozłączne
części widma: punktową (brak iniektywności $T-zI$),
ciągłą (iniektywność i gęsty, lecz niepełny obraz)
oraz residualną (iniektywność i obraz, który nie jest
gęsty). Wyczerpują one widmo, ponieważ bijekcja
$T-zI$ ma na mocy powyższego argumentu ograniczoną
odwrotność.
\end{definicja}

\begin{uwaga}[Tożsamość rezolwenty i otwartość zbioru rezolwenty]
Jeżeli $z\in\rho(T)$ i $|w-z|\|R_z\|<1$, to
na $D(T)$ mamy faktoryzację
$T-wI=(I-(w-z)R_z)(T-zI)$.
Szereg Neumanna dla operatora $(w-z)R_z$ zapewnia
odwracalność pierwszego czynnika w $\mathcal L(H)$.
Drugi czynnik jest bijekcją $D(T)\to H$.
Stąd $T-wI$ jest bijekcją i jego ograniczona odwrotność
ma postać $R_w=R_z(I-(w-z)R_z)^{-1}$, z obrazem
w $D(T)$. Zatem $w\in\rho(T)$. W szczególności
$\rho(T)$ jest otwarty, a $\sigma(T)$ domknięte.
Dla dowolnych $z,w\in\rho(T)$ mamy
\[
 R_z-R_w=R_z[(T-wI)-(T-zI)]R_w
 =(z-w)R_zR_w.
\]
Rachunek jest poprawny dziedzinowo, ponieważ
$R_w(H)\subset D(T)$, a $R_z(T-zI)$ jest identycznością
na $D(T)$. Zamieniając $z$ z $w$, otrzymujemy
także $(z-w)R_wR_z=R_z-R_w$, więc rezolwenty komutują
(dla $z=w$ jest to bezpośrednie).
Jeżeli $R_z$ jest zwarty dla jednego $z$,
to $R_w=R_z+(w-z)R_zR_w$ jest zwarty dla każdego
$w\in\rho(T)$: złożenie zwartego $R_z$ z ograniczonym
$R_w$ jest zwarte, a suma operatorów zwartych
jest zwarta.
\end{uwaga}

\begin{twierdzenie}[Alternatywa Fredholma dla laplasjanu]
\label{thm:pde12-fredholm}
Niech $A=-\Delta_D$ będzie operatorem na obszarze ograniczonym Lipschitzowskim $\Omega$, a $\zeta\in\C$.
Jeżeli $\zeta\notin\{\lambda_j:j\ge1\}$, to dla
każdego $f\in L^2(\Omega;\C)$ istnieje dokładnie jedno
$u\in D(A)$ spełniające $(A-\zeta I)u=f$, i
\[
\|u\|_2\le\frac{\|f\|_2}
{\operatorname{dist}(\zeta,\{\lambda_j\})}.
\]
Jeżeli $\zeta$ jest wartością własną, to równanie
$(A-\zeta I)u=f$ jest rozwiązywalne wtedy i tylko wtedy,
gdy $f\perp\ker(A-\zeta I)$; rozwiązanie jest wówczas
jednoznaczne z dokładnością do dodania elementu jądra.
Rezolwenta $(A-\zeta I)^{-1}$ w każdym punkcie
$\zeta\in\rho(A)$ jest zwarta, a
$\sigma(A)=\{\lambda_j:j\ge1\}$ jako zbiór (bez krotności).
\end{twierdzenie}
\begin{proof}
Przechodzimy do zespolenia przestrzeni i operatora:
$D(A_\C)=D(A)+iD(A)$ oraz
$A_\C(v+iw)=Av+iAw$. Zachowujemy oznaczenie $A$.
Rzeczywista baza własna pozostaje bazą ortonormalną
po zespoleniu, gdyż rozwijamy oddzielnie część rzeczywistą
i urojoną. Charakterystyki z twierdzenia
~\ref{thm:pde12-formdomain} obowiązują z modułami
współczynników zespolonych przez dodanie odpowiednich
tożsamości dla obu części. W szczególności operator
jest domknięty także po zespoleniu (stosujemy kryterium
ciągowe osobno do obu części).
Operator zespolony jest także samosprzężony.
Istotnie, jest symetryczny przez rozwinięcie zespolonego
iloczynu skalarnego i rzeczywistą symetrię $A$.
Jeśli $v+iw\in D(A_\C^*)$, $v,w\in H$ rzeczywistym,
i $A_\C^*(v+iw)=p+iq$, to testowanie definicji
sprzężenia rzeczywistymi $h\in D(A)$ daje
$(Ah,v)=(h,p)$ oraz $(Ah,w)=(h,q)$.
Rzeczywista samosprzężoność implikuje $v,w\in D(A)$
oraz $Av=p$, $Aw=q$. Wraz z symetrią daje to równość
dziedzin i operatorów.
Piszemy $f=\sum_jf_j\phi_j$, a dla potencjalnego
rozwiązania $u=\sum_ju_j\phi_j$, gdzie
$f_j=(f,\phi_j)_H$, $u_j=(u,\phi_j)_H$.
Z charakterystyki $D(A)$ wynika, że równanie jest
równoważne warunkom $(\lambda_j-\zeta)u_j=f_j$
dla wszystkich $j$, o ile $u\in D(A)$.
Jeżeli $\zeta$ nie jest wartością własną, to odległość
$\delta=\inf_j|\lambda_j-\zeta|$ jest dodatnia:
$\lambda_j\to\infty$, więc istnieje $N$ takie, że
$|\lambda_j-\zeta|>|\lambda_1-\zeta|+1$ dla $j>N$.
Minimum odległości jest zatem osiągane wśród pierwszych
$N$ wyrazów i jest dodatnie. Definiujemy $u_j=f_j/(\lambda_j-\zeta)$.
Z Parsevala
$\|u\|_2^2=\sum_j|f_j|^2/|\lambda_j-\zeta|^2
\le\delta^{-2}\|f\|_2^2$.
Dodatkowo z nierówności trójkąta
$\lambda_j\le|\lambda_j-\zeta|+|\zeta|$ otrzymujemy
jednostajne oszacowanie
$\lambda_j/|\lambda_j-\zeta|\le1+|\zeta|/\delta$.
W szczególności $\sum_j|u_j|^2<\infty$ pozwala
zdefiniować $u$ jako granicę szeregu ortogonalnego. Stąd $\sum\lambda_j^2|u_j|^2<\infty$,
więc $u\in D(A)$ na mocy~\ref{thm:pde12-formdomain},
i rzeczywiście rozwiązuje równanie. Każde inne rozwiązanie
ma te same współczynniki, co dowodzi jednoznaczności.
Oszacowanie normy rezolwenty jest ostre: wybierając
$f=\phi_{j_0}$ dla indeksu realizującego $\delta$,
otrzymujemy $\|f\|_2=1$ oraz $\|R_\zeta f\|_2=\delta^{-1}$.

Jeśli $\zeta$ jest wartością własną, to dla indeksów
$J=\{j:\lambda_j=\zeta\}$ równanie daje warunek
$f_j=0$ dla każdego $j\in J$.
Z charakterystyki spektralnej dziedziny $A$ wynika
$\ker(A-\zeta I)=\operatorname{span}\{\phi_j:j\in J\}$:
wektory z jądra mają zerowe współczynniki poza $J$,
a każda taka suma skończona należy do $D(A)$ i spełnia
równanie własne. Warunek na $f_j$ jest zatem dokładnie
ortogonalnością $f$ do tego jądra.
Dla $j\notin J$ możemy określić
$u_j=f_j/(\lambda_j-\zeta)$, a współczynniki $u_j$
z $j\in J$ pozostają dowolne. Zbiór $J$ jest skończony
na mocy skończonej krotności, a jego dopełnienie
jest nieskończone. Na tym dopełnieniu
$\delta'=\min_{j\notin J}|\lambda_j-\zeta|>0$
z tego samego argumentu co poprzednio, rozpoczętego
od dowolnego indeksu spoza $J$.
Mamy
\[
 \sum_{j\notin J}\lambda_j^2|u_j|^2
 \le(1+|\zeta|/\delta')^2\sum_{j\notin J}|f_j|^2.
\]
Pozostałe składniki są skończoną sumą, więc
otrzymujemy $u\in D(A)$ i równanie dla wszystkich
współczynników. Wybór $u_j=0$ dla $j\in J$ daje jedyne
rozwiązanie $u_0$ ortogonalne do jądra, z oszacowaniem
$\|u_0\|_2\le(\delta')^{-1}\|f\|_2$.
Wszystkie rozwiązania mają postać $u_0+w$,
$w\in\ker(A-\zeta I)$; z ortogonalności
$\|u_0+w\|_2^2=\|u_0\|_2^2+\|w\|_2^2$.
Zatem $u_0$ jest jedynym rozwiązaniem o najmniejszej
normie $L^2$.

Dla $\zeta$ poza zbiorem wartości własnych definiujemy
na całym $H$ operator skończonego rzędu
$R_{\zeta,N}f=\sum_{j\le N} f_j(\lambda_j-\zeta)^{-1}\phi_j$.
Z Parsevala
\[
 \|(R_\zeta-R_{\zeta,N})f\|_2^2
 \le \sup_{j>N}|\lambda_j-\zeta|^{-2}\|f\|_2^2.
\]
Supremum dąży do zera, bo $\lambda_j\to\infty$.
Każdy $R_{\zeta,N}$ jest zwarty jako ograniczony operator
skończonego rzędu. Dla dowolnego $\varepsilon>0$ wybieramy
$N$ tak, aby norma różnicy była mniejsza niż $\varepsilon/2$.
Obraz kuli jednostkowej przez $R_{\zeta,N}$ ma skończoną
$(\varepsilon/2)$-sieć, która jest $\varepsilon$-siecią
obrazu przez $R_\zeta$. Z całkowitej ograniczoności
i zupełności $H$ wynika zwartość $R_\zeta$.
Każda $\zeta$ spoza zbioru wartości własnych należy do
$\rho(A)$, a każda wartość własna należy do $\sigma(A)$
z powodu braku iniektywności. To daje równość zbiorów.
\end{proof}

Ogólny mechanizm alternatywy Fredholma dla operatorów
zwartych omówiono również w \atlasref{FA-30::thm:fredholm-hilbert30}{FA-30: alternatywa Fredholma dla operatora zwartego}.

\begin{twierdzenie}[Alternatywa Fredholma dla operatorów zwartych]
\label{thm:pde12-general-fredholm}
Niech $H$ będzie zespoloną przestrzenią Hilberta,
a $K:H\to H$ operatorem zwartym. Wówczas
$T=I-K$ ma domknięty obraz oraz skończeniewymiarowe
jądra $\ker T$ i $\ker T^*$. Zachodzi
$\dim\ker T=\dim\ker T^*$, a dla $f\in H$
równanie $(I-K)u=f$ ma rozwiązanie wtedy i tylko wtedy,
gdy $f\perp\ker(I-K^*)$. Rozwiązanie jest jednoznaczne
z dokładnością do dodania elementu $\ker(I-K)$.
\end{twierdzenie}
\begin{proof}
\textbf{Krok 1: skończony wymiar jądra.}
Jeśli $\ker T$ byłoby nieskończeniewymiarowe,
wybralibyśmy w nim nieskończony układ ortonormalny
$(e_n)$. Równość $Te_n=0$ oznacza $Ke_n=e_n$.
Ponieważ $(e_n)$ nie ma podciągu zbieżnego, byłoby
to sprzeczne ze zwartością $K$. Skończoną wymiarowość
$\ker T^*$ otrzymamy w kroku 4.

\textbf{Krok 2: domkniętość obrazu.} Niech
$N=\ker T$ i $M=N^\perp$. Jądro jest domknięte
jako jądro ciągłego operatora, więc $H=N\oplus M$.
Pokażemy, że istnieje $c>0$
takie, że $\|Tx\|\ge c\|x\|$ dla wszystkich $x\in M$.
Dla $M=\{0\}$ można wybrać dowolne $c>0$.
W przeciwnym razie, gdyby oszacowanie było fałszywe, można wybrać $x_n\in M$,
$\|x_n\|=1$, $Tx_n\to0$. Zwarty operator $K$
daje podciąg $(Kx_{n_k})$ zbieżny w $H$.
Ponieważ $x_{n_k}=Tx_{n_k}+Kx_{n_k}$,
ten sam podciąg $x_{n_k}$ zbiega do pewnego $x\in M$
o normie $1$. Ciągłość $T$ daje $Tx=0$,
czyli $x\in M\cap N=\{0\}$; sprzeczność.
Teraz niech $Tx_n\to y$ w $H$.
Zastępując $x_n$ przez jego rzut na $M$, nie zmieniamy
$Tx_n$. Uzyskane $x_n\in M$ są Cauchy'ego,
bo $\|x_n-x_m\|\le c^{-1}\|Tx_n-Tx_m\|$.
Zatem $x_n\to x\in M$ i $Tx=y$, dowodząc domkniętości
$\operatorname{Ran}T$.

\textbf{Krok 3: równość wymiarów jąder.}
Pokażemy najpierw, że $K$ można aproksymować
w normie operatorowej operatorami skończonego rzędu.
Dla $0<\varepsilon<1$ zbiór $K(B_H)$ jest względnie
zwarty, więc można znaleźć skończony zbiór
$y_1,\dots,y_m$ w $H$ taki, że każdy $Kx$, $\|x\|\le1$,
leży w odległości mniejszej niż $\varepsilon$ od
jednego z $y_j$. Jeśli $W=\operatorname{span}\{y_j\}$,
a $P_W$ oznacza rzut ortogonalny na $W$, to
$F=P_WK$ jest skończonego rzędu i
$\|K-F\|\le\varepsilon$, gdyż
$\|Kx-P_WKx\|=\operatorname{dist}(Kx,W)
\le\min_j\|Kx-y_j\|<\varepsilon$ dla $\|x\|\le1$.
Zapiszmy $K=F+E$, gdzie $\|E\|\le\varepsilon<1$.
Szereg $\sum_{n=0}^\infty E^n$ zbiega w normie operatorowej
z oszacowania $\|E^n\|\le\|E\|^n$.
Mnożąc sumy częściowe z obu stron przez $I-E$
i przechodząc do granicy, otrzymujemy odwrotność $I-E$,
a faktoryzacja
\[
I-K=(I-E)(I-G),\qquad G=(I-E)^{-1}F
\]
przedstawia $T$ jako iloczyn operatora odwracalnego
i operatora $I-G$ z $G$ skończonego rzędu.
Niech $W'=\operatorname{Ran}G$ (przestrzeń skończeniewymiarowa).
W rozkładzie $H=W'\oplus(W')^\perp$ operator $I-G$
ma postać blokową
\[
(I-G)(w,z)=((I-G|_{W'})w-Gz,\ z).
\]
Przez $I_{W'}$ oznaczamy identyczność na $W'$.
Druga współrzędna wektora z jądra musi być zerowa,
więc jądro jest równe
$\ker(I_{W'}-G|_{W'})\oplus\{0\}$.
Oznaczmy $S=I_{W'}-G|_{W'}$ i $L=\operatorname{Ran}S$.
Wprost z postaci blokowej $(a,z)$ należy do obrazu
$I-G$ wtedy i tylko wtedy, gdy $a+Gz\in L$.
Odwzorowanie $(a,z)\mapsto a+Gz+L$ do przestrzeni
ilorazowej $W'/L$ jest liniowe i surjektywne
(bierzemy $z=0$); jego jądrem jest właśnie
$\operatorname{Ran}(I-G)$. Z twierdzenia o izomorfizmie
przestrzeń ilorazowa $H/\operatorname{Ran}(I-G)$
jest izomorficzna z $W'/L$, więc kowymiar obrazu
wynosi $\dim W'-\dim L$.
Z twierdzenia o wymiarze jądra i obrazu w $W'$
ten kowymiar równa się $\dim\ker S$.
Mnożenie z lewej przez odwracalny $I-E$ nie zmienia
jądra, a na przestrzeniach ilorazowych indukuje izomorfizm
$H/\operatorname{Ran}(I-G)\to H/\operatorname{Ran}T$,
$[x]\mapsto[(I-E)x]$, z odwrotnością indukowaną przez
$(I-E)^{-1}$. Nie zmienia zatem kowymiaru obrazu. Otrzymujemy
$\dim\ker T=\operatorname{codim}\operatorname{Ran}T$.

\textbf{Krok 4: warunek rozwiązywalności.}
Dla dowolnego ograniczonego operatora w Hilbercie
$(\operatorname{Ran}T)^\perp=\ker T^*$,
gdyż $(Tu,v)=0$ dla wszystkich $u$ jest dokładnie
warunkiem $T^*v=0$. Obraz $T$ jest domknięty
z kroku 2, stąd
$\operatorname{Ran}T=(\ker T^*)^\perp$.
Ponieważ $\operatorname{Ran}T$ jest domknięty,
rozkład ortogonalny
$H=\operatorname{Ran}T\oplus(\operatorname{Ran}T)^\perp$
identyfikuje przestrzeń ilorazową
$H/\operatorname{Ran}T$ z $(\operatorname{Ran}T)^\perp$
przez rzut ortogonalny.
Stąd $\dim\ker T^*=\operatorname{codim}\operatorname{Ran}T
=\dim\ker T$ na mocy kroku 3. Dowodzi to także
warunku rozwiązywalności dla $f$.
Każde dwa rozwiązania różnią się o element jądra,
a dodanie dowolnego jego elementu do rozwiązania zachowuje
równanie. Jeśli warunek zgodności jest spełniony, rozkładamy
dowolne rozwiązanie $u=u_N+u_M$ w $N\oplus M$.
Ponieważ $Tu_N=0$, wektor $u_M$ również jest rozwiązaniem.
Dwa rozwiązania w $M$ różniłyby się o wektor
z $M\cap N=\{0\}$, więc rozwiązanie w $M$ jest jedyne.
Z kroku 2 ma ono normę co najwyżej $c^{-1}\|f\|$;
jest też rozwiązaniem o najmniejszej normie,
gdyż $\|u_M+n\|^2=\|u_M\|^2+\|n\|^2$ dla $n\in N$.
W szczególności jeśli $\ker T=\{0\}$, równość wymiarów
daje $\ker T^*=\{0\}$, a więc $T$ jest bijekcją
z ograniczoną odwrotnością. Jeśli $T$ jest surjektywny,
to $\ker T^*=\{0\}$, a równość wymiarów daje iniektywność.
\end{proof}

\begin{uwaga}
W odróżnieniu od twierdzenia~\ref{thm:pde12-fredholm},
nie założyliśmy tutaj samosprzężoności $K$.
Dlatego w warunku rozwiązywalności pojawia się
$\ker(I-K^*)$, a niekoniecznie $\ker(I-K)$.
W terminologii teorii operatorów wynik stwierdza,
że $I-K$ jest operatorem Fredholma o indeksie $0$,
czyli $\dim\ker T-\operatorname{codim}\operatorname{Ran}T=0$.
Na przykład na $\C^2$ dla
$K=\begin{pmatrix}1&-1\\0&1\end{pmatrix}$
mamy $T=I-K=\begin{pmatrix}0&1\\0&0\end{pmatrix}$.
W standardowej bazie ortonormalnej $e_1,e_2$ mamy
$T(u_1,u_2)=(u_2,0)$ oraz $T^*(v_1,v_2)=(0,v_1)$.
Zatem $\ker T=\C e_1$, $\ker T^*=\C e_2$,
a równanie $Tu=f$ jest rozwiązywalne dokładnie dla
$f=(f_1,0)$ (wówczas $u_2=f_1$, $u_1$ dowolne).
Na przykład $f=e_2$ jest ortogonalny do $\ker T$,
ale równanie jest nierozwiązywalne, natomiast $f=e_1$
jest rozwiązywalny mimo braku ortogonalności do $\ker T$.
To pokazuje, dlaczego sprzężenie w warunku zgodności
jest niezbędne.
\end{uwaga}


\begin{kontrprzyklad}[Brak zwartości na całej przestrzeni]
\label{ex:pde12-unbounded-domain}
W $H=L^2(\R^d;\C)$ rozważmy $A=-\Delta$ z dziedziną
$H^2(\R^d;\C)$. Na mocy twierdzenia
~\ref{thm:pde8-plancherel} transformata
$U=(2\pi)^{-d/2}\mathcal F$ jest unitarna w konwencji
rozdziału 8. Z reguł różniczkowania dystrybucji
(twierdzenie~\ref{thm:pde8-dist-rules}) wynika
$UAU^{-1}=M_m$, gdzie $m(\xi)=|\xi|^2$ i
$D(M_m)=\{g\in L^2:mg\in L^2\}$.
Istotnie, jeśli $u\in H^2$, to $\partial_j^2u\in L^2$
dla każdego $j$, a w sensie dystrybucji temperowanych
$U(\partial_j^2u)=-\xi_j^2Uu$.
Zatem $|\xi|^2Uu=-\sum_{j=1}^d U(\partial_j^2u)\in L^2$.
Odwrotnie, jeśli $g$ i $|\xi|^2 g$ należą do $L^2$,
to wszystkie iloczyny $(i\xi)^\alpha g$ dla
$|\alpha|\le2$ należą do $L^2$, ponieważ
$|\xi^\alpha|\le 1+|\xi|^2$. Odwrotna transformata
tych iloczynów daje słabe pochodne $U^{-1}g$.

Operator $M_m$ jest gęsto określony, gdyż zawiera
wszystkie funkcje $L^2$ o ograniczonym nośniku, a ich
gęstość wynika z obcięcia do kul i zbieżności w $L^2$.
Jest symetryczny, ponieważ $m$ jest rzeczywista.
Jeśli $h\in D(M_m^*)$ oraz $k=M_m^*h$, to dla każdego
$g\in L^2$ o ograniczonym nośniku
$\int mg\overline h\,\mathrm d\lambda_d
=\int g\overline k\,\mathrm d\lambda_d$.
Ponieważ $m$ jest ograniczona na każdej kuli,
$g=\chi_{B(0,R)}(mh-k)$ należy do $L^2$
i do $D(M_m)$. Wstawiając ten test, otrzymujemy
$\|\chi_{B(0,R)}(mh-k)\|_2^2=0$ dla każdego $R>0$.
Stąd $mh=k\in L^2$, czyli $h\in D(M_m)$.
Wraz z symetrią daje to samosprzężoność $M_m$.
Ponieważ $U$ jest unitarny, definicja sprzężenia daje
$(U^{-1}M_mU)^*=U^{-1}M_m^*U$, z dziedziną
$U^{-1}D(M_m^*)$. Zatem $A$ również jest samosprzężony,
a oba operatory są domknięte.

Dla $z\notin[0,\infty)$ mnożnik $(m-z)^{-1}$ jest
ograniczony przez $\operatorname{dist}(z,[0,\infty))^{-1}$,
a $m/(m-z)=1+z/(m-z)$ także jest ograniczony.
Mnożenie przez $(m-z)^{-1}$ jest więc ograniczoną
odwrotnością $M_m-zI$ o obrazie w $D(M_m)$.
Dla $s\ge0$ zbiór
$E_n=\{\xi:\big||\xi|^2-s\big|<1/n\}$ ma dodatnią
skończoną miarę: jest otwarty i niepusty (zawiera
punkt o normie $\sqrt{s}$) oraz zawiera się w kuli
o promieniu $\sqrt{s+1/n}$. Funkcje $g_n=\chi_{E_n}/\sqrt{\lambda_d(E_n)}$
należą do $D(M_m)$, mają normę $1$ i spełniają
$\|(M_m-sI)g_n\|_2\le1/n$. Wyklucza to istnienie
ograniczonej odwrotności. Jednocześnie zbiór
$\{\xi:|\xi|^2=s\}$ ma miarę zero (dla $s=0$ jest
punktem, dla $s>0$ sferą; miara sfery jest zerowa
przez oszacowanie miarą pierścieni o promieniach
dążących z obu stron do $\sqrt{s}$).
Jeśli $(M_m-sI)g=0$, to $g$ znika prawie wszędzie
poza tym zbiorem, a więc $g=0$ w $L^2$.
Dopełnienie ortogonalne obrazu to jądro sprzężenia
$(M_m-sI)^*=M_m-sI$, z definicji sprzężenia
i rzeczywistości $s$. Obraz jest zatem gęsty,
ale nie jest całym $H$: gdyby był całym $H$,
operator byłby bijekcją i miałby ograniczoną odwrotność,
co przeczy ciągowi $g_n$.
Z definicji każdy $s\ge0$ należy do widma ciągłego.
Unitarna równoważność przenosi iniektywność, gęstość
i surjektywność obrazu oraz ograniczoność odwrotności,
więc $\sigma(A)=[0,\infty)$ jest w całości widmem ciągłym.

Pokażmy niezależnie brak zwartości rezolwenty $B=(A+I)^{-1}$.
Wybierzmy $0\ne f\in C_c^\infty(\R^d)$ oraz $L>0$
tak duże, aby przesunięcia
$f_n(x)=f(x-nLe_1)$ miały parami rozłączne nośniki.
Mamy $\|f_n\|_2=\|f\|_2$, zatem ciąg jest ograniczony.
Operator $B=U^{-1}M_{(1+m)^{-1}}U$ komutuje
z translacjami: transformata przesunięcia o $a$
jest mnożeniem przez $e^{-ia\cdot\xi}$, co wynika
ze zmiany zmiennych dla funkcji całkowalnych
i przez gęstość oraz unitarność dla $L^2$.
Stąd $u_n=Bf_n=u(\cdot-nLe_1)$, gdzie $u=Bf\ne0$
dzięki iniektywności $B$.

Dla dowolnego $v\in C_c^\infty(\R^d)$ o nośniku $Q$
nierówność Cauchy'ego--Schwarza daje
\[
 |(u_n,v)|\le\|v\|_2
 \left(\int_{Q-nLe_1}|u(y)|^2\,\mathrm d\lambda_d(y)
 \right)^{1/2}\longrightarrow0,
\]
gdyż zbiory $Q-nLe_1$ opuszczają każdą ustaloną kulę,
a $|u|^2$ jest całkowalna. Dla ogólnego $v\in L^2$
przybliżamy go testem $v_\varepsilon$ z błędem
$\|v-v_\varepsilon\|_2<\varepsilon$;
$|(u_n,v-v_\varepsilon)|\le\|u\|_2\varepsilon$
jednostajnie względem $n$. Otrzymujemy $u_n\rightharpoonup0$.
Ponieważ $\|u_n\|_2=\|u\|_2>0$, żaden podciąg nie może
zbiegać silnie: jego granica musiałaby być jednocześnie
granicą słabą $0$. Zwartość $B$ przeczyłaby temu faktowi.
Nie należy więc oczekiwać dyskretnego rozwinięcia
po funkcjach własnych jak na obszarze ograniczonym.
\end{kontrprzyklad}

\section{Półgrupa cieplna zbudowana z widma}
\label{sec:pde12-heat}

\begin{definicja}[Silnie ciągła półgrupa]
Niech $H$ będzie przestrzenią Banacha. Rodzinę
$S(t)\in\mathcal L(H)$, $t\ge0$, nazywamy
\emph{silnie ciągłą półgrupą}, gdy $S(0)=I$,
$S(t+s)=S(t)S(s)$ dla wszystkich $s,t\ge0$, a dla
każdego $u\in H$ zachodzi $\|S(t)u-u\|_H\to0$
przy $t\downarrow0$. Jej \emph{generator} $G$ jest
dany przez $Gu=\lim_{t\downarrow0}(S(t)u-u)/t$
na dziedzinie tych $u$, dla których granica istnieje w $H$.
Dziedzina $D(G)$ jest podprzestrzenią liniową, a $G$ jest liniowy,
gdyż ilorazy różnicowe są liniowe i granica w normie
zachowuje kombinacje liniowe.
\end{definicja}

\begin{twierdzenie}[Półgrupa generowana przez $-A$]
\label{thm:pde12-heat}
Niech $A=-\Delta_D$ dla obszaru ograniczonego Lipschitzowskiego $\Omega$, a $H=L^2(\Omega)$.
Dla $t\ge0$ określmy
\[
S(t)u=\sum_{j=1}^\infty e^{-\lambda_jt}(u,\phi_j)_H\phi_j.
\]
Wtedy $(S(t))_{t\ge0}$ jest silnie ciągłą półgrupą na $H$
o generatorze dokładnie $-A$ i
$\|S(t)\|_{H\to H}=e^{-\lambda_1t}$.
Dla $t>0$, całkowitego $m\ge1$ i $u\in H$ zachodzi
$S(t)u\in D(A^m)$ oraz
\[
\|A^mS(t)u\|_H
\le\bigg(\frac{m}{e t}\bigg)^m\|u\|_H.
\]
Dla $u_0\in D(A)$ funkcja $u(t)=S(t)u_0$
należy do $C^1([0,\infty);H)\cap C([0,\infty);D(A))$
(z normą grafową na $D(A)$) i spełnia $u'(t)+Au(t)=0$, $u(0)=u_0$.
Dla dowolnego $u_0\in H$ to samo równanie zachodzi dla
$t>0$, a $u(t)\to u_0$ w $H$ przy $t\downarrow0$.
\end{twierdzenie}
\begin{proof}
\textbf{Definicja i własność półgrupy.}
Niech $u_j=(u,\phi_j)_H$. Ponieważ
$\sum_j e^{-2\lambda_jt}|u_j|^2\le\sum_j|u_j|^2<\infty$,
sumy częściowe są ciągiem Cauchy'ego w $H$ dzięki
ortonormalności. Ich granica definiuje liniowy operator.
Z Parsevala
$\|S(t)u\|_H^2=\sum_je^{-2\lambda_jt}|u_j|^2
\le e^{-2\lambda_1t}\|u\|_H^2$, więc jest on ograniczony. Dla $u=\phi_1$ mamy
równość, co daje dokładną normę. Mnożenie współczynników
$e^{-\lambda_js}e^{-\lambda_jt}=e^{-\lambda_j(s+t)}$
pokazuje $S(t)S(s)=S(t+s)$, a $S(0)=I$.

\textbf{Ciągłość silna.} Dla $u\in H$ mamy
$\|S(t)u-u\|_H^2=\sum_j|e^{-\lambda_jt}-1|^2|u_j|^2$.
Dla każdego $j$ czynnik w module dąży do zera,
a jest ograniczony przez $1$. Ponieważ
$\sum_j|u_j|^2<\infty$, z twierdzenia o zbieżności
zmajoryzowanej dla szeregu otrzymujemy zbieżność do zera.
Dla $h>0$ mamy
$\|S(t+h)u-S(t)u\|\le\|S(h)u-u\|$,
a dla $0<h\le t$ również
$\|S(t)u-S(t-h)u\|\le\|S(h)u-u\|$.
W obu nierównościach zastosowaliśmy własność półgrupy
i $\|S(r)\|\le1$ dla $r\ge0$. Zatem $S(\cdot)u$
jest ciągła w każdym $t\ge0$.

\textbf{Wygładzanie.} Najpierw uzasadnijmy charakterystykę
potęg operatora:
\[
 D(A^m)=\{u\in H:\sum_j\lambda_j^{2m}|u_j|^2<\infty\},
 \qquad A^mu=\sum_j\lambda_j^m u_j\phi_j .
\]
Dla $m=1$ jest to twierdzenie~\ref{thm:pde12-formdomain}.
Jeżeli wzór zachodzi dla $m$, to z definicji
$u\in D(A^{m+1})$ dokładnie wtedy, gdy $u\in D(A^m)$
i $A^mu\in D(A)$. Jeśli $u\in D(A^{m+1})$, z założenia indukcyjnego
współczynniki $A^mu$ wynoszą $\lambda_j^mu_j$,
a warunek $A^mu\in D(A)$ daje
$\sum_j\lambda_j^{2m+2}|u_j|^2<\infty$.
Odwrotnie, jeśli ta suma jest skończona, to również
$\sum_j\lambda_j^{2m}|u_j|^2<\infty$, gdyż
$\lambda_j^{-2}\le\lambda_1^{-2}$. Najpierw otrzymujemy
zatem $u\in D(A^m)$, a następnie $A^mu\in D(A)$
z charakterystyki $D(A)$, zastosowanej do współczynników
$\lambda_j^mu_j$. Zastosowanie $A$ do $A^mu$
daje wzór dla $A^{m+1}$.
Dla $m\ge1$ i $t>0$
$\sup_{x\ge0}x^me^{-tx}=(m/(et))^m$:
pochodna dla $x>0$ wynosi
$x^{m-1}e^{-tx}(m-tx)$, więc funkcja rośnie do $m/t$
i maleje dalej, a w zerze ma wartość $0$.
Wobec tego
$\sum_j\lambda_j^{2m}e^{-2\lambda_jt}|u_j|^2
\le(m/(et))^{2m}\|u\|_H^2$.
Z wykazanej charakterystyki dziedzin potęg wynika
$S(t)u\in D(A^m)$, a Parseval daje podane oszacowanie.
Dokładna norma wynosi
$\|A^mS(t)\|=\sup_j\lambda_j^me^{-t\lambda_j}$:
nierówność w jedną stronę wynika z Parsevala, a w drugą
z podstawienia kolejnych jednostkowych $\phi_j$.
Oszacowanie przez maksimum po wszystkich $x\ge0$
nie musi być równością dla dyskretnego widma.

\textbf{Generator.} Dla $u\in D(A)$ iloraz
$(e^{-\lambda_jt}-1)/t$ dąży do $-\lambda_j$, a
$|(e^{-\lambda_jt}-1)/t|\le\lambda_j$,
ponieważ $1-e^{-a}=\int_0^a e^{-r}\,\mathrm d\lambda_1(r)
\le a$ dla $a\ge0$. Kwadrat modułu różnicy
między współczynnikiem ilorazu a $-\lambda_j u_j$
jest ograniczony przez $4\lambda_j^2|u_j|^2$.
Sumowalność tej majoranty pozwala zastosować zbieżność
zmajoryzowaną i otrzymać
$\lim_{t\downarrow0}(S(t)u-u)/t=-Au$ w $H$.
Odwrotnie, gdy granica ilorazu istnieje w $H$ i równa się
$g$, jej $j$-ty współczynnik musi równać się
$-\lambda_ju_j$. Stąd z Parsevala
$\sum_j\lambda_j^2|u_j|^2=\|g\|_H^2<\infty$,
a więc $u\in D(A)$. To dowodzi równości dziedzin generatora.
Dla $u_0\in D(A)$ charakterystyka spektralna daje
$S(t)u_0\in D(A)$ oraz $AS(t)u_0=S(t)Au_0$.
Ciągłość w normie grafowej wynika więc z silnej ciągłości
$S(\cdot)$ zastosowanej do $u_0$ i $Au_0$.
Prawostronny iloraz różnicowy w chwili $t$ jest równy
$S(t)(S(h)-I)u_0/h$ i dąży do $-S(t)Au_0$.
Dla $t>0$ lewostronny iloraz jest równy
$S(t-h)(S(h)-I)u_0/h$ (jest to iloraz
$(u(t)-u(t-h))/h$ dla $0<h<t$), a jego różnica względem
$-S(t)Au_0$ jest ograniczona przez
\[
 \|(S(h)-I)u_0/h+Au_0\|
 +\|(S(t)-S(t-h))Au_0\|\longrightarrow0.
\]
Pochodna $-S(t)Au_0$ jest ciągła również w zerze
(w zerze rozumiemy pochodną prawostronną).
Dla dowolnego $u_0\in H$ i $t_0>0$ możemy zapisać
$u(t)=S(t-t_0/2)S(t_0/2)u_0$ dla $t\ge t_0/2$.
Wygładzanie daje $S(t_0/2)u_0\in D(A)$,
więc udowodniony wynik stosuje się w otoczeniu $t_0$.
W ten sposób otrzymujemy równanie i ciągłą pochodną
dla wszystkich $t>0$, a w zerze silną ciągłość.
Dla $t>0$ można iterować ten argument: dla każdego
całkowitego $k\ge1$ mamy $u^{(k)}(t)=(-A)^kS(t)u_0$,
ponieważ na otoczeniu dowolnego $t_0>0$ początkowy
wektor $S(t_0/2)u_0$ należy do wszystkich $D(A^k)$.
Indukcyjnie różniczkujemy $S(t-t_0/2)A^kS(t_0/2)u_0$,
korzystając z wykazanej komutacji i przynależności
wektora po prawej stronie do $D(A)$.
Jednoznaczność w klasie rozwiązań ciągłych w $H$
od $t=0$ i klasycznych dla $t>0$ wynika z równań
skalarnych: dla $v'(t)+Av(t)=0$ współczynniki
$v_j(t)=(v(t),\phi_j)_H$ spełniają
$v_j'+\lambda_jv_j=0$. Zatem
$e^{\lambda_jt}v_j(t)=v_j(0)$ przez ciągłość w zerze,
a zupełność bazy daje $v(t)=S(t)v(0)$.
\end{proof}

\begin{twierdzenie}[Wzór Duhamela w $L^2$]
\label{thm:pde12-duhamel}
Niech $T>0$, $u_0\in H=L^2(\Omega)$ i
$f\in C([0,T];H)$, gdzie $S(t)$ jest półgrupą
z twierdzenia~\ref{thm:pde12-heat}. Funkcja
\[
u(t)=S(t)u_0+\int_0^tS(t-s)f(s)\,\mathrm d\lambda_1(s),
\qquad 0\le t\le T,
\]
gdzie całkę względem miary Lebesgue'a na osi czasu
rozumiemy w sensie Bochnera, jest ciągła
z wartościami w $H$ i spełnia oszacowanie
$\|u(t)\|_H\le e^{-\lambda_1t}\|u_0\|_H
+\int_0^te^{-\lambda_1(t-s)}\|f(s)\|_H\,\mathrm d\lambda_1(s)$.
Jest jedynym rozwiązaniem \emph{łagodnym} zagadnienia
$u'(t)+Au(t)=f(t)$, $u(0)=u_0$, zdefiniowanym
przez powyższy wzór całkowy. Jeśli dodatkowo
$u_0\in D(A)$ i $f\in C^1([0,T];H)$, to
$u\in C^1([0,T];H)\cap C([0,T];D(A))$
(z normą grafową na $D(A)$) i równanie zachodzi w $H$
dla każdego $t\in[0,T]$, z pochodnymi jednostronnymi
w punktach końcowych.
\end{twierdzenie}
\begin{proof}
Mapowanie $(t,s)\mapsto S(t-s)f(s)$ jest silnie ciągłe
na trójkącie $0\le s\le t\le T$: dla $(t_n,s_n)\to(t,s)$
różnica wartości jest ograniczona przez
$\|f(s_n)-f(s)\|+\|(S(t_n-s_n)-S(t-s))f(s)\|$,
co dąży do zera. Ciągła funkcja o wartościach w przestrzeni
Banacha na przedziale zwartym jest całkowalna w sensie
Bochnera. Korzystamy tu z kryterium całkowalności
i oszacowania normy całki z analizy funkcjonalnej
(\atlasref{FA-15B::thm:bochner-criterion15b}{kryterium
całkowalności Bochnera} oraz
\atlasref{FA-15B::prop:riemann-bochner15b}{zgodność
całek Riemanna i Bochnera dla funkcji ciągłych}). Stosując dokładną normę $S(t-s)$, dostajemy
podane oszacowanie.
Niech $M=\max_{s\in[0,T]}\|f(s)\|$. Dla $t_n\to t$
na wspólnym przedziale $[0,\min(t_n,t)]$ różnica
całkowanych funkcji dąży jednostajnie do zera
z jednostajnej ciągłości na zwartym trójkącie.
Norma całki na pozostałym odcinku jest nie większa
niż $M|t_n-t|$. To dowodzi ciągłości składnika całkowego;
w szczególności jego wartość w zerze wynosi $0$.

Dla silniejszych danych, po zmianie zmiennej $r=t-s$,
możemy pisać $F(t)=\int_0^tS(r)f(t-r)\,\mathrm d\lambda_1(r)$.
Dla $h>0$ i $t+h\le T$ rozkładamy iloraz różnicowy:
\[
 \begin{aligned}
 \frac{F(t+h)-F(t)}h
 &=\int_0^t S(r)\frac{f(t+h-r)-f(t-r)}h
 \,\mathrm d\lambda_1(r)\\
 &\quad+\frac1h\int_t^{t+h}S(r)f(t+h-r)
 \,\mathrm d\lambda_1(r).
 \end{aligned}
\]
Na pierwszym przedziale ilorazy
$(f(t+h-r)-f(t-r))/h$ zbiegają jednostajnie
do $f'(t-r)$ z jednostajnej ciągłości $f'$ i wzoru
Newtona--Leibniza w przestrzeni Banacha.
Mnożenie przez $S(r)$ nie powiększa błędów.
Średnia całkowa na krótkim odcinku zbiega do $S(t)f(0)$
przez wspólną ciągłość $S(r)f(t+h-r)$.
Dla ilorazu lewostronnego, przy $0<h\le t$, piszemy
\[
 \begin{aligned}
 \frac{F(t)-F(t-h)}h
 &=\int_0^{t-h} S(r)\frac{f(t-r)-f(t-h-r)}h
 \,\mathrm d\lambda_1(r)\\
 &\quad+\frac1h\int_{t-h}^tS(r)f(t-r)
 \,\mathrm d\lambda_1(r).
 \end{aligned}
\]
Na pierwszym przedziale błąd względem całki z $S(r)f'(t-r)$
dąży do zera jednostajnie, a norma brakującej całki po
$[t-h,t]$ nie przekracza $h\max\|f'\|$. Drugi składnik
dąży do $S(t)f(0)$ przez wspólną ciągłość. Otrzymujemy
więc tę samą granicę.
Stąd
$F'(t)=S(t)f(0)+\int_0^tS(r)f'(t-r)\,\mathrm d\lambda_1(r)$,
z pochodnymi jednostronnymi na końcach przedziału.
Pochodna jest ciągła przez ten sam argument co dla $F$,
z $f'$ zamiast $f$.
Ustalmy $t>0$ i $0<\varepsilon<t$.
Na $[\varepsilon,t]$ funkcja $r\mapsto S(r)f(t-r)$
jest ciągła w normie grafowej: zapisujemy
$S(r)=S(\varepsilon/2)S(r-\varepsilon/2)$, gdzie
$S(\varepsilon/2):H\to D(A)$ jest ograniczony dzięki
oszacowaniu wygładzania dla $m=1$. Dziedzina $D(A)$
z normą grafową jest przestrzenią Banacha, więc całka
$F_\varepsilon(t)=\int_\varepsilon^t S(r)f(t-r)\,
\mathrm d\lambda_1(r)$ istnieje w tej przestrzeni.
Odwzorowania $D(A)\to H$, $v\mapsto v$ i $v\mapsto Av$
są ograniczone w normie grafowej, zatem można je
przenieść pod całkę Bochnera (np. przez granicę sum
Riemanna dla funkcji ciągłej).
Z poprzedniego twierdzenia i reguły iloczynu dla
silnej pochodnej mamy dla $r\in[\varepsilon,t]$
\[
 \frac{\mathrm d}{\mathrm dr}[S(r)f(t-r)]
 =-AS(r)f(t-r)-S(r)f'(t-r).
\]
Reguła iloczynu wynika z rozdzielenia ilorazu różnicowego
na zmianę $S(r)$ przy ustalonym wektorze i zmianę
wektora $f(t-r)$; stosujemy $\|S(r)\|\le1$ oraz silną
ciągłość. Pochodna jest ciągła na $[\varepsilon,t]$:
pierwszy składnik jest ciągły dzięki wykazanej ciągłości
grafowej, drugi dzięki wspólnej ciągłości $S(r)f'(t-r)$.
Można więc zastosować wzór Newtona--Leibniza w $H$.
Całkując tę tożsamość, dostajemy
\[
 AF_\varepsilon(t)=S(\varepsilon)f(t-\varepsilon)
 -S(t)f(0)-\int_\varepsilon^t S(r)f'(t-r)\,
 \mathrm d\lambda_1(r).
\]
Gdy $\varepsilon\downarrow0$, mamy
$F_\varepsilon(t)\to F(t)$ w $H$, a prawa strona
zbiega w $H$ do $f(t)-F'(t)$: pierwszy składnik
zbiega przez wspólną ciągłość, a pominięta całka
ma normę co najwyżej $\varepsilon\max\|f'\|$.
Domkniętość $A$ daje $F(t)\in D(A)$ i
$AF(t)=f(t)-F'(t)$. W zerze $F(0)=0$ oraz
$F'(0)=f(0)$, więc tożsamość zachodzi także tam.
Ponieważ $F$ i $AF=f-F'$ są ciągłe w $H$,
$F$ jest ciągła w normie grafowej.
Dla $S(t)u_0$ mamy $\frac{\mathrm d}{\mathrm dt}S(t)u_0
=-AS(t)u_0$ z poprzedniego twierdzenia.
Zsumowanie obu rachunków daje $u'+Au=f$.

Dla dowolnego $u_0\in H$ i ciągłego $f$ wynik definiuje
rozwiązanie łagodne. Jednoznaczność w tej klasie wynika
bezpośrednio z wymogu spełniania wzoru całkowego.
Jeśli $v$ jest rozwiązaniem klasycznym, to
dla każdego $j$ funkcja $v_j(t)=(v(t),\phi_j)_H$
spełnia $v_j'(t)+\lambda_j v_j(t)=f_j(t)$, gdzie
$f_j(t)=(f(t),\phi_j)_H$. Pomnożenie przez
$e^{\lambda_jt}$ i całkowanie daje
\[
 v_j(t)=e^{-\lambda_jt}(u_0,\phi_j)_H+
 \int_0^t e^{-\lambda_j(t-s)} f_j(s)\,
 \mathrm d\lambda_1(s).
\]
Prawa strona jest $j$-tym współczynnikiem wzoru Duhamela,
ponieważ ograniczony funkcjonał $(\cdot,\phi_j)_H$
przenosi się pod całkę Bochnera. Z zupełności bazy
otrzymujemy $v=u$, czyli rozwiązanie łagodne zgadza się
z klasycznym, ilekroć drugie istnieje. W szczególności
sama ciągłość $f$ nie wystarcza na ogół do
przynależności rozwiązania do $D(A)$ w każdej chwili,
jak pokazuje poniższy kontrprzykład.
\end{proof}

\begin{kontrprzyklad}[Ciągła prawa strona bez regularności w dziedzinie operatora]
Ustalmy $T=1$ i $u_0=0$. Wybierzmy rosnący ciąg indeksów
$j_n$ taki, że $L_n=\lambda_{j_n}>4$ oraz
$L_{n+1}\ge4L_n$; jest to możliwe, gdyż $\lambda_j\to\infty$.
Niech $\eta:\R\to[0,1]$ będzie ciągła, niezerowa i z nośnikiem
w $[1,2]$, np. $\eta(r)=\max(0,1-2|r-3/2|)$.
Dla $0\le s\le1$ określmy
\[
 f(s)=\sum_{n\ge1}n^{-1/2}\eta(L_n(1-s))\phi_{j_n}.
\]
Niezerowe składniki mogą występować tylko na przedziałach
$[1-2/L_n,1-1/L_n]$, które są parami rozłączne.
W każdej chwili co najwyżej jeden składnik jest niezerowy,
a przedziały gromadzą się tylko przy $s=1$.
Funkcja jest ciągła na $[0,1)$ dzięki ciągłości $\eta$,
a w $s=1$ również, ponieważ $f(1)=0$ i amplitudy
$n^{-1/2}$ dążą do zera. Zatem $f\in C([0,1];H)$.
W chwili $t=1$ współczynniki rozwiązania łagodnego wynoszą
\[
 \begin{aligned}
 (u(1),\phi_{j_n})_H
 &=n^{-1/2}\int_0^1 e^{-L_n(1-s)}
 \eta(L_n(1-s))\,\mathrm d\lambda_1(s)\\
 &=\frac{C_\eta}{\sqrt n\,L_n},
 \qquad
 C_\eta=\int_1^2 e^{-r}\eta(r)\,\mathrm d\lambda_1(r)>0 .
 \end{aligned}
\]
Pozostałe współczynniki znikają. Druga równość wynika
z podstawienia $r=L_n(1-s)$ i nośnika $\eta$.
Zatem $\sum_j\lambda_j^2|(u(1),\phi_j)_H|^2
=C_\eta^2\sum_{n\ge1}1/n=\infty$.
Twierdzenie~\ref{thm:pde12-formdomain} daje
$u(1)\notin D(A)$ mimo ciągłości $f$ i zerowego $u_0$.
\end{kontrprzyklad}

\begin{uwaga}[Regularność w czasie i na brzegu]
Oszacowanie $\|A^mS(t)\|\lesssim t^{-m}$ jest
oszacowaniem operatorowym. Nie daje automatycznie
$C^\infty(\overline\Omega)$ na dowolnym obszarze Lipschitzowskim: w narożach może zabraknąć regularności
brzegowej. Dla jednorodnego równania ciepła i $t>0$
$u(t)=S(t)u_0$ jest jednak gładka przestrzennie we wnętrzu.
Istotnie, dla każdego $m\ge1$ wektory $u(t),Au(t),\ldots,
A^{m-1}u(t)$ należą do $D(A)\subset H_0^1$,
a $A^mu(t)\in L^2$. Stosując wewnętrzną regularność
eliptyczną (twierdzenia~\ref{thm:pde15-h2}
i~\ref{thm:pde15-bootstrap}) kolejno do
$A^{m-1}u(t),\ldots,u(t)$, otrzymujemy
$A^{m-1}u(t)\in H^2_{\mathrm{loc}}$,
$A^{m-2}u(t)\in H^4_{\mathrm{loc}}$ i w końcu
$u(t)\in H^{2m}_{\mathrm{loc}}$.
Ponieważ $m$ jest dowolne, dla każdego zwartego
podzbioru wnętrza i dowolnego $k\ge0$ wybieramy
$2m>k+d/2$. Lokalne zanurzenia Sobolewa dają
reprezentanta klasy $C^k$ w otoczeniu tego zbioru.
Reprezentanci dla różnych $k$ i otoczeń zgadzają się
na częściach wspólnych, gdyż są funkcjami ciągłymi
zgodnymi prawie wszędzie. Otrzymujemy reprezentanta klasy $C^\infty(\Omega)$.
Ten argument nie dotyczy dowolnej niegładkiej prawej strony.
\end{uwaga}

\section{Równanie falowe: propagacja i zachowanie energii}
\label{sec:pde12-wave}

Piszemy $V=H_0^1(\Omega)$, $H=L^2(\Omega)$ i utożsamiamy
$H$ z podprzestrzenią $V^*=H^{-1}(\Omega)$ za pomocą
odwzorowania $h\mapsto(v\mapsto(h,v)_H)$.
Jest ono ciągłe z nierówności Poincarégo i różnowartościowe
z gęstości $V$ w $H$. W tej sekcji pracujemy w przestrzeniach
rzeczywistych; wyniki dla danych zespolonych otrzymujemy
przez rozdzielenie części rzeczywistej i urojonej.
Rozszerzenie energetyczne operatora $A$ oznaczamy przez
$\mathcal A:V\to V^*$, gdzie
\[
 \langle\mathcal Au,v\rangle
 =\int_\Omega\nabla u(x)\cdot\nabla v(x)\,\mathrm d\lambda_d(x).
\]
Jest ono ograniczone z nierówności Cauchy'ego--Schwarza
i zgodne z $A$ na $D(A)$.
Przez rozwiązanie energetyczne rozumiemy funkcję
$u\in C(\R;V)\cap C^1(\R;H)$ spełniającą dane początkowe
oraz dla każdego $v\in V$ i $\psi\in C_c^\infty(\R)$
tożsamość
\[
 \int_\R (u(t),v)_H\psi''(t)\,\mathrm d\lambda_1(t)
 +c^2\int_\R\langle\mathcal Au(t),v\rangle\psi(t)\,
 \mathrm d\lambda_1(t)=0.
\]
W tym sensie $u_{tt}+c^2\mathcal Au=0$ jest równaniem
o wartościach w $V^*$; nie wymagamy $u(t)\in D(A)$.

\begin{twierdzenie}[Spektralne rozwiązanie równania falowego]
\label{thm:pde12-wave}
Niech $A=-\Delta_D$ jak wyżej, $c>0$,
$u_0\in H_0^1(\Omega)$ i $u_1\in L^2(\Omega)$.
Oznaczmy $a_j=(u_0,\phi_j)$, $b_j=(u_1,\phi_j)$.
Wtedy szereg
\[
u(t)=\sum_{j=1}^\infty
\left(a_j\cos(c\sqrt{\lambda_j}t)
+\frac{b_j}{c\sqrt{\lambda_j}}\sin(c\sqrt{\lambda_j}t)\right)\phi_j
\]
wyznacza jedyne rozwiązanie energetyczne równania
$u_{tt}+c^2Au=0$ z danymi $u(0)=u_0$, $u_t(0)=u_1$,
przy czym
$u\in C(\R;H_0^1(\Omega))\cap C^1(\R;L^2(\Omega))$.
Energia
\[
E(t)=\frac12\|u_t(t)\|_2^2+
\frac{c^2}{2}\|\nabla u(t)\|_2^2
\]
jest stała i równa $E(0)$.
\end{twierdzenie}
\begin{proof}
Z twierdzenia~\ref{thm:pde12-formdomain} mamy
$\sum_j\lambda_j|a_j|^2<\infty$ i
$\sum_j|b_j|^2<\infty$. Dla współczynnika
$u_j(t)=a_j\cos(c\sqrt{\lambda_j}t)
+b_j\sin(c\sqrt{\lambda_j}t)/(c\sqrt{\lambda_j})$
zachodzi
\[
|u_j'(t)|^2+c^2\lambda_j|u_j(t)|^2
=|b_j|^2+c^2\lambda_j|a_j|^2.
\]
Istotnie, dla $\omega_j=c\sqrt{\lambda_j}>0$
para $(\omega_j u_j(t),u_j'(t))$ jest równa
\[
 \begin{pmatrix}\cos(\omega_jt)&\sin(\omega_jt)\\
 -\sin(\omega_jt)&\cos(\omega_jt)\end{pmatrix}
 \begin{pmatrix}\omega_j a_j\\b_j\end{pmatrix}.
\]
Macierz jest ortogonalna, więc zachowuje sumę kwadratów
modułów obu współrzędnych (również zespolonych).
Oznaczmy przez $u_N$ sumę pierwszych $N$ modów.
Dla $M>N$, z Parsevala i charakterystyki dziedziny formy,
\begin{align*}
 \|u_M'(t)-u_N'(t)\|_2^2
 &+c^2\|\nabla(u_M(t)-u_N(t))\|_2^2\\
 &=\sum_{j=N+1}^{M}(|b_j|^2+c^2\lambda_j|a_j|^2).
\end{align*}
Prawa strona nie zależy od $t$ i dąży do zera
jednostajnie względem $M>N$. Ponieważ norma gradientu
jest równoważna pełnej normie $V$, z zupełności
$C([-T,T];V)$ i $C([-T,T];H)$ uzyskujemy jednostajne
granice $u_N\to u$ i $u_N'\to p$ na każdym takim przedziale.
Granice są zgodne na częściach wspólnych przedziałów.
W tożsamości
$u_N(t)-u_N(0)=\int_0^t u_N'(s)\,\mathrm d\lambda_1(s)$
przechodzimy do granicy w $H$ (dla $t<0$ używamy
całki zorientowanej). Z ciągłości $p$ i twierdzenia
Newtona--Leibniza wynika $u'=p$ w $H$.
Dane początkowe otrzymujemy z rozwinięć $u_N(0)$ w $V$
oraz $u_N'(0)$ w $H$.

Każde $u_N$ jest gładkie w czasie i spełnia
$u_N''+c^2Au_N=0$ w $H$, ponieważ każdy współczynnik
spełnia skalarne równanie oscylatora.
Dla $v\in V$ mnożymy równanie przez $\psi$ i całkujemy
dwukrotnie przez części w czasie. W otrzymanej tożsamości
błędy przy zastąpieniu $u_N$ przez $u$ są ograniczone
przez
$\sup_{|t|\le T}\|u_N(t)-u(t)\|_H\|v\|_H\|\psi''\|_{L^1}$
oraz
$c^2\sup_{|t|\le T}\|u_N(t)-u(t)\|_V\|v\|_V\|\psi\|_{L^1}$,
gdy nośnik $\psi$ leży w $[-T,T]$.
Oba błędy znikają, co daje tożsamość definiującą
rozwiązanie energetyczne.

Dla jednoznaczności niech $w$ będzie różnicą dwóch
rozwiązań energetycznych z zerowymi danymi. Wybierając
$v=\phi_j$ w tożsamości słabej, z równania własnego
otrzymujemy $w_j''=-c^2\lambda_j w_j$ dystrybucyjnie,
gdzie $w_j=(w,\phi_j)_H\in C^1(\R)$.
Prawa strona jest ciągła, więc $w_j'$ jest klasycznie
różniczkowalna: różnica między $w_j'$ a funkcją
$t\mapsto\int_0^t(-c^2\lambda_jw_j(s))\,\mathrm d\lambda_1(s)$
ma zerową pochodną dystrybucyjną, a zatem jest stała.
W konsekwencji $w_j\in C^2(\R)$ i spełnia klasyczne
równanie oscylatora z $w_j(0)=w_j'(0)=0$.
Jednoznaczność tego równania daje $w_j=0$ dla każdego $j$.
Zupełność bazy daje $w=0$.
Dla sum częściowych energia jest stała dzięki tożsamości
modalnej. Silna zbieżność obu składników energii oraz
danych początkowych daje po przejściu do granicy
$E(t)=E(0)$ dla każdego $t\in\R$.
\end{proof}

\begin{uwaga}[Silniejsze dane początkowe]
Jeśli ponadto $u_0\in D(A)$ i $u_1\in V$, to
$u\in C(\R;D(A))\cap C^1(\R;V)\cap C^2(\R;H)$
i równanie zachodzi w $H$. Istotnie, mnożymy tożsamość
energii każdego modu przez $\lambda_j$.
Sumowalność $\sum_j(\lambda_j|b_j|^2+
c^2\lambda_j^2|a_j|^2)$ daje jednostajną zbieżność
ogonów $u_N$ w normie grafowej $A$ oraz $u_N'$ w $V$;
normę $L^2$ kontroluje $\lambda_1>0$.
Ponadto $u_N''=-c^2Au_N$ zbiega jednostajnie w $H$
na każdym przedziale zwartym. Dwukrotne zastosowanie
argumentu z całkowaniem po czasie dowodzi podanej
różniczkowalności i równania.
\end{uwaga}

\begin{uwaga}[Czego analiza spektralna nie zastępuje?]
Rozwinięcie w funkcje własne pokazuje istnienie i
zachowanie globalnej energii, ale nie ujawnia bezpośrednio
geometrii stożków charakterystycznych. Skończona prędkość
propagacji, uzyskana w twierdzeniu~\ref{thm:pde5-cone}
metodą energii lokalnej (dla rozważanych tam założeń),
pozostaje niezależną, silniejszą lokalnie informacją.
\end{uwaga}

\section{Przepływ Schrödingera i związek z operatorem unitarnym}
\label{sec:pde12-schroedinger}

\begin{twierdzenie}[Ewolucja unitarna laplasjanu Dirichleta]
\label{thm:pde12-unitary}
Dla laplasjanu Dirichleta $A$ na ograniczonym obszarze
Lipschitzowskim $\Omega$ przechodzimy do zespolonej przestrzeni
$H=L^2(\Omega;\C)$ i zespolonego operatora $A$.
Zdefiniujmy $U(t)u=\sum_je^{-it\lambda_j}(u,\phi_j)_H\phi_j$,
$t\in\R$. Wtedy $(U(t))_{t\in\R}$ jest silnie
ciągłą grupą unitarną z generatorem $-iA$.
Dla każdego $u_0\in H$ zachodzi $\|U(t)u_0\|_2=\|u_0\|_2$.
Jeśli $u_0\in D(A)$, to $u(t)=U(t)u_0$ jest jedynym
rozwiązaniem $iu_t=Au$, $u(0)=u_0$ w klasie
$C^1(\R;H)\cap C(\R;D(A))$, gdzie $D(A)$ ma normę grafową.
\end{twierdzenie}
\begin{proof}
Piszemy $u_j=(u,\phi_j)_H$. Ponieważ $|e^{-it\lambda_j}|=1$
i $\sum_j|u_j|^2<\infty$, szereg jest zbieżny w $H$
z kryterium Cauchy'ego dla szeregu ortogonalnego.
Z liniowości granicy otrzymujemy operator liniowy,
a Parseval daje $\|U(t)u\|_2=\|u\|_2$.
Na współczynnikach odczytujemy
$U(t)U(s)=U(t+s)$ oraz $U(0)=I$.
Zatem $U(t)$ jest surjektywną izometrią liniową
z odwrotnością $U(-t)$, a więc operatorem unitarnym
(z tożsamości polaryzacyjnej).
Silna ciągłość w zerze wynika z
$\sum_j|e^{-it\lambda_j}-1|^2|u_j|^2\to0$
przez zbieżność zmajoryzowaną, z majorantą
$4|u_j|^2$. Ciągłość w pozostałych punktach wynika
z własności grupy i izometrii.

Dla $u\in D(A)$ i $h\ne0$ mamy
\[
 \left|\frac{e^{-ih\lambda_j}-1}{h}\right|
 =\left|\frac{1}{h}\int_0^h(-i\lambda_j)e^{-ir\lambda_j}
 \,\mathrm d\lambda_1(r)\right|\le\lambda_j .
\]
Iloraz dąży do $-i\lambda_j$ dla $h\to0$.
Kwadrat modułu różnicy odpowiednich współczynników
jest ograniczony przez $4\lambda_j^2|u_j|^2$,
co jest sumowalne. Zatem
$\lim_{h\to0}(U(h)u-u)/h=-iAu$ w $H$.
Odwrotnie, jeśli iloraz ma granicę $g\in H$,
ciągłość funkcjonałów $(\cdot,\phi_j)_H$ daje
$(g,\phi_j)_H=-i\lambda_j u_j$ dla każdego $j$.
Parseval daje $\sum_j\lambda_j^2|u_j|^2=\|g\|_H^2<\infty$,
czyli $u\in D(A)$.
Generator jest więc dokładnie $-iA$.

Z charakterystyki spektralnej wynika, że $U(t)$ zachowuje
$D(A)$ i $AU(t)u=U(t)Au$ na $D(A)$.
Ciągłość grafowa wynika z silnej ciągłości zastosowanej
do $u$ i $Au$. Iloraz różnicowy w dowolnej chwili $t$
jest równy $U(t)(U(h)u-u)/h$; stąd pochodna
$-iU(t)Au$ istnieje i jest ciągła w $H$.
To dowodzi równania Schrödingera. Dla innego rozwiązania
$v$ w podanej klasie współczynniki spełniają
$v_j'=-i\lambda_jv_j$, zatem
$v_j(t)=e^{-it\lambda_j}(u_0,\phi_j)_H$.
Zupełność bazy dowodzi jednoznaczności.
Ponadto dla każdego $u\in V$ tożsamość spektralna energii
daje $\|\nabla U(t)u\|_2=\|\nabla u\|_2$; dla $u\in D(A)$
zachodzi $\|AU(t)u\|_2=\|Au\|_2$.
Przepływ zachowuje więc również normę formy i normę grafową.
\end{proof}

\begin{uwaga}
Dla ogólnego operatora samosprzężonego, niekoniecznie
mającego zwartą rezolwentę, rachunek spektralny daje
silnie ciągłą grupę unitarną $e^{-itA}$.
Półgrupa $e^{-tA}$ składa się z operatorów ograniczonych
dla wszystkich $t\ge0$, jeśli $A$ jest ograniczony z dołu;
jeżeli $\sigma(A)\subset[a,\infty)$, to
$\|e^{-tA}\|\le e^{-ta}$. Bez dolnego ograniczenia widma
nie można w ogólności definiować w ten sposób półgrupy
ograniczonych operatorów na całym $H$.
Twierdzenie Stone'a identyfikuje generator każdej silnie
ciągłej grupy unitarnej jako $-iA$ dla dokładnie jednego
operatora samosprzężonego $A$ (konwencja $e^{-itA}$;
dla $e^{itA}$ generator ma znak przeciwny).
Ogólny dowód i rachunek spektralny podano w analizie
funkcjonalnej, rozdział 33
(\atlasref{FA-33::thm:stone33}{twierdzenie Stone'a}).
Tutaj dowiedliśmy bezpośrednio konstrukcji dla laplasjanu
Dirichleta z dyskretnym widmem.
\end{uwaga}

\section{Podsumowanie: co wynika z założeń, a co nie?}
\label{sec:pde12-summary}

Na obszarze ograniczonym Lipschitzowskim
słaba forma Dirichleta definiuje samosprzężony
operator $A=-\Delta_D$. Koercywność formy przesuniętej zapewnia ciągłość
$(A+I)^{-1}:L^2\to H_0^1$, a zwartość zanurzenia
$H_0^1\hookrightarrow L^2$ zapewnia zwartość rezolwenty.
W połączeniu z samosprzężonością i różnowartościowością
rezolwenty otrzymujemy dyskretne widmo i ortonormalną
bazę funkcji własnych. Stąd otrzymujemy alternatywę
Fredholma, rozwinięcia dla równania ciepła i falowego
oraz przepływ unitarny równania Schrödingera.

Nie dowiedliśmy tu ogólnego twierdzenia spektralnego dla
nieograniczonych operatorów samosprzężonych ani
pełnego twierdzenia Hille'a--Yosidy. W kolejnych rozdziałach
wykorzystamy jednak już udowodnione przypadki oraz
wyraźnie oddzielimy twierdzenia w pełnej ogólności od
wyników obowiązujących dla zwartej rezolwenty.

% Rozdział 13. Tekst przeznaczony do włączenia przed \appendix.
